\clearpage
\section{Copying and parametricity: checked intermediate declarations}
\label{sec:remaining-template-concordance}
These are supporting results rather than completed template proofs. Copying retains every original input and resource distributor. The parametricity prefix uses the actual private store and arbitrary stable argument. Universal residual synthesis and robustness remain open as specified in Appendix B.1.

\subsection{Bound Weakening Splice}
Exact beta geometry for an ignored bound argument. The argument is retained, with every original output, attached to the actual weakening.

The complete unused argument, including its original resource outputs.
\lean{Copying/BoundWeakeningSplice}{13}{13}{BoundWeakeningSplice.discarded}
\par\penalty0

This defines the left graph fragment for the construction just described. Boundary positions are kept distinct from newly inserted internal vertices.
\lean{Copying/BoundWeakeningSplice}{15}{16}{BoundWeakeningSplice.left}
\par\penalty0

This defines the right graph fragment for the construction just described. Boundary positions are kept distinct from newly inserted internal vertices.
\lean{Copying/BoundWeakeningSplice}{18}{19}{BoundWeakeningSplice.right}
\par\penalty0

This defines the internal-vertex map for the construction just described. Boundary positions are kept distinct from newly inserted internal vertices.
\lean{Copying/BoundWeakeningSplice}{21}{44}{BoundWeakeningSplice.vertices}
\par\penalty0

This defines the edge map for the construction just described. Boundary positions are kept distinct from newly inserted internal vertices.
\lean{Copying/BoundWeakeningSplice}{46}{58}{BoundWeakeningSplice.edges}
\par\penalty0

This defines the left Body Port data for the construction just described. Boundary positions are kept distinct from newly inserted internal vertices.
\lean{Copying/BoundWeakeningSplice}{60}{61}{BoundWeakeningSplice.leftBodyPort}
\par\penalty0

This defines the right Body Port data for the construction just described. Boundary positions are kept distinct from newly inserted internal vertices.
\lean{Copying/BoundWeakeningSplice}{63}{65}{BoundWeakeningSplice.rightBodyPort}
\par\penalty0

This defines the left Argument Port data for the construction just described. Boundary positions are kept distinct from newly inserted internal vertices.
\lean{Copying/BoundWeakeningSplice}{67}{68}{BoundWeakeningSplice.leftArgumentPort}
\par\penalty0

This defines the right Argument Port data for the construction just described. Boundary positions are kept distinct from newly inserted internal vertices.
\lean{Copying/BoundWeakeningSplice}{70}{73}{BoundWeakeningSplice.rightArgumentPort}
\par\penalty0

The body Port commutes assertion is the compatibility law for the indicated component of this construction. It relates the actual finite coordinates or certificates, with the hypotheses stated in the linked declaration.
For the proof, constructor and boundary cases are separated explicitly; the defined maps and incidence equations are substituted.
\lean{Copying/BoundWeakeningSplice}{75}{85}{BoundWeakeningSplice.bodyPort_commutes}
\par\penalty0

The argument Port commutes assertion is the compatibility law for the indicated component of this construction. It relates the actual finite coordinates or certificates, with the hypotheses stated in the linked declaration.
For the proof, constructor and boundary cases are separated explicitly; the defined maps and incidence equations are substituted.
\lean{Copying/BoundWeakeningSplice}{105}{116}{BoundWeakeningSplice.argumentPort_commutes}
\par\penalty0

This defines the left Body Edge data for the construction just described. Boundary positions are kept distinct from newly inserted internal vertices.
\lean{Copying/BoundWeakeningSplice}{118}{118}{BoundWeakeningSplice.leftBodyEdge}
\par\penalty0

This defines the right Body Edge data for the construction just described. Boundary positions are kept distinct from newly inserted internal vertices.
\lean{Copying/BoundWeakeningSplice}{119}{119}{BoundWeakeningSplice.rightBodyEdge}
\par\penalty0

This defines the left Argument Edge data for the construction just described. Boundary positions are kept distinct from newly inserted internal vertices.
\lean{Copying/BoundWeakeningSplice}{120}{120}{BoundWeakeningSplice.leftArgumentEdge}
\par\penalty0

This defines the right Argument Edge data for the construction just described. Boundary positions are kept distinct from newly inserted internal vertices.
\lean{Copying/BoundWeakeningSplice}{121}{121}{BoundWeakeningSplice.rightArgumentEdge}
\par\penalty0

This defines the left Weakening Edge data for the construction just described. Boundary positions are kept distinct from newly inserted internal vertices.
\lean{Copying/BoundWeakeningSplice}{122}{123}{BoundWeakeningSplice.leftWeakeningEdge}
\par\penalty0

This defines the right Weakening Edge data for the construction just described. Boundary positions are kept distinct from newly inserted internal vertices.
\lean{Copying/BoundWeakeningSplice}{124}{125}{BoundWeakeningSplice.rightWeakeningEdge}
\par\penalty0

The left body source identity gives the entire ordered source list after applying the construction, rather than only its underlying set.
For the proof, constructor and boundary cases are separated explicitly; the defined maps and incidence equations are substituted.
\lean{Copying/BoundWeakeningSplice}{127}{133}{BoundWeakeningSplice.left_body_source}
\par\penalty0

The left body target identity gives the entire ordered target list after applying the construction, with output order retained.
For the proof, constructor and boundary cases are separated explicitly; the defined maps and incidence equations are substituted.
\lean{Copying/BoundWeakeningSplice}{135}{141}{BoundWeakeningSplice.left_body_target}
\par\penalty0

The right body source identity gives the entire ordered source list after applying the construction, rather than only its underlying set.
For the proof, constructor and boundary cases are separated explicitly.
\lean{Copying/BoundWeakeningSplice}{143}{145}{BoundWeakeningSplice.right_body_source}
\par\penalty0

The right body target identity gives the entire ordered target list after applying the construction, with output order retained.
For the proof, constructor and boundary cases are separated explicitly.
\lean{Copying/BoundWeakeningSplice}{147}{149}{BoundWeakeningSplice.right_body_target}
\par\penalty0

The left argument source identity gives the entire ordered source list after applying the construction, rather than only its underlying set.
For the proof, constructor and boundary cases are separated explicitly.
\lean{Copying/BoundWeakeningSplice}{151}{153}{BoundWeakeningSplice.left_argument_source}
\par\penalty0

The left argument target identity gives the entire ordered target list after applying the construction, with output order retained.
For the proof, constructor and boundary cases are separated explicitly.
\lean{Copying/BoundWeakeningSplice}{155}{157}{BoundWeakeningSplice.left_argument_target}
\par\penalty0

The right argument source identity gives the entire ordered source list after applying the construction, rather than only its underlying set.
For the proof, constructor and boundary cases are separated explicitly.
\lean{Copying/BoundWeakeningSplice}{159}{161}{BoundWeakeningSplice.right_argument_source}
\par\penalty0

The right argument target identity gives the entire ordered target list after applying the construction, with output order retained.
For the proof, constructor and boundary cases are separated explicitly.
\lean{Copying/BoundWeakeningSplice}{163}{165}{BoundWeakeningSplice.right_argument_target}
\par\penalty0

The weakening target identity gives the entire ordered target list after applying the construction, with output order retained.
For the proof, constructor and boundary cases are separated explicitly; the defined maps and incidence equations are substituted.
\lean{Copying/BoundWeakeningSplice}{167}{177}{BoundWeakeningSplice.weakening_target}
\par\penalty0

The edge cases assertion is the compatibility law for the indicated component of this construction. It relates the actual finite coordinates or certificates, with the hypotheses stated in the linked declaration.
For the proof, constructor and boundary cases are separated explicitly; the defined maps and incidence equations are substituted.
\lean{Copying/BoundWeakeningSplice}{179}{187}{BoundWeakeningSplice.edge_cases}
\par\penalty0

The edges body assertion is the compatibility law for the indicated component of this construction. It relates the actual finite coordinates or certificates, with the hypotheses stated in the linked declaration.
For the proof, constructor and boundary cases are separated explicitly; the defined maps and incidence equations are substituted.
\lean{Copying/BoundWeakeningSplice}{189}{190}{BoundWeakeningSplice.edges_body}
\par\penalty0

The edges argument assertion is the compatibility law for the indicated component of this construction. It relates the actual finite coordinates or certificates, with the hypotheses stated in the linked declaration.
For the proof, constructor and boundary cases are separated explicitly; the defined maps and incidence equations are substituted.
\lean{Copying/BoundWeakeningSplice}{191}{192}{BoundWeakeningSplice.edges_argument}
\par\penalty0

The edges weakening assertion is the compatibility law for the indicated component of this construction. It relates the actual finite coordinates or certificates, with the hypotheses stated in the linked declaration.
For the proof, constructor and boundary cases are separated explicitly; the defined maps and incidence equations are substituted.
\lean{Copying/BoundWeakeningSplice}{193}{194}{BoundWeakeningSplice.edges_weakening}
\par\penalty0

Beta does not erase an ignored stable argument. It exposes precisely a separate weakening connected to that entire argument, preserving all edges, recursive box labels, output order and the active input.
\lean{Copying/BoundWeakeningSplice}{199}{228}{BoundWeakeningSplice.iso}
\par\penalty0

\subsection{Copying After Input Permutation}
Input permutations fixing the selected input commute with the original copy reduct. Both comparisons are actual copies, so proved successor coherence supplies a recursive graph isomorphism without modifying any local rule.

Reindex the complete original source, keeping every payload edge and ordered output, before invoking either actual copy transition.
\lean{Copying/CopyingAfterInputPermutation}{14}{20}{CopyingInputPermutation.beforeReorderedIso}
\par\penalty0

For every original copyable payload, a permutation fixing input zero preserves the literal copy reduct up to recursive graph representation.
\lean{Copying/CopyingAfterInputPermutation}{24}{36}{CopyingInputPermutation.afterIso}
\par\penalty0

\subsection{Copying Double}
Literal two-copy endpoints. The second selected original input is copied inside the retained branch of the first actual Figure 17 reduct. Both original two-use distributors and every residual weakening remain in these graphs.

Replace the retained branch of the literal two-branch copying construction. The active payload and the full ordered term/atom output permutation stay fixed.
\lean{Copying/CopyingDouble}{14}{19}{CopyingDouble.copiesWith}
\par\penalty0

The original outer resource distributors follow the chosen retained branch.
\lean{Copying/CopyingDouble}{22}{25}{CopyingDouble.afterWith}
\par\penalty0

Two literal copies, first at original input zero and then at original input one. The latter is input zero of the first residual contraction tree.
\lean{Copying/CopyingDouble}{29}{32}{CopyingDouble.afterSecond}
\par\penalty0

Before the second step, this retained-branch construction is exactly the first literal reduct, with no sharing-tree normalization.
Both sides are definitionally equal.
\lean{Copying/CopyingDouble}{36}{38}{CopyingDouble.beforeSecond}
\par\penalty0

Exchange the two selected original inputs, preserving every other position.
\lean{Copying/CopyingDouble}{41}{42}{CopyingDouble.exchangeInputs}
\par\penalty0

Perform the same two literal copies in the opposite order and restore the original external input interface. This is the exact interchange endpoint.
\lean{Copying/CopyingDouble}{46}{51}{CopyingDouble.opposite}
\par\penalty0

\subsection{Copying Double At}
Normalize an arbitrary second selected input to position one, perform the literal two-copy diagram, and restore the complete original input ordering.

This defines the ordering data for the construction just described. Boundary positions are kept distinct from newly inserted internal vertices.
\lean{Copying/CopyingDoubleAt}{15}{16}{CopyingDoubleAt.ordering}
\par\penalty0

The ordering one assertion is the compatibility law for the indicated component of this construction. It relates the actual finite coordinates or certificates, with the hypotheses stated in the linked declaration.
For the proof, the defined maps and incidence equations are substituted.
\lean{Copying/CopyingDoubleAt}{18}{19}{CopyingDoubleAt.ordering_one}
\par\penalty0

The ordering zero assertion is the compatibility law for the indicated component of this construction. It relates the actual finite coordinates or certificates, with the hypotheses stated in the linked declaration.
For the proof, the remaining natural-number index bounds follow by linear arithmetic; the defined maps and incidence equations are substituted.
\lean{Copying/CopyingDoubleAt}{21}{27}{CopyingDoubleAt.ordering_zero}
\par\penalty0

This defines the normalized data for the construction just described. Boundary positions are kept distinct from newly inserted internal vertices.
\lean{Copying/CopyingDoubleAt}{29}{29}{CopyingDoubleAt.normalized}
\par\penalty0

The actual first-copy successor when the entered original input is j.
\lean{Copying/CopyingDoubleAt}{32}{34}{CopyingDoubleAt.leftAfter}
\par\penalty0

The actual second-copy successor on the original template's right side.
\lean{Copying/CopyingDoubleAt}{37}{38}{CopyingDoubleAt.rightAfter}
\par\penalty0

Source equality here is an isomorphism of full ordered graphs, obtained from two actual original copies and proved successor coherence.
\lean{Copying/CopyingDoubleAt}{42}{45}{CopyingDoubleAt.rightSourceIso}
\par\penalty0

The first original copy at input one, in the normalized presentation.
For the proof, the defined maps and incidence equations are substituted.
\lean{Copying/CopyingDoubleAt}{48}{58}{CopyingDoubleAt.normalized_left_rewrite}
\par\penalty0

No branch or input is omitted: this is an actual original rewrite at j.
For the proof, the defined maps and incidence equations are substituted.
\lean{Copying/CopyingDoubleAt}{61}{68}{CopyingDoubleAt.left_rewrite}
\par\penalty0

The other side performs its next actual original copy at the same j.
For the proof, the defined maps and incidence equations are substituted.
\lean{Copying/CopyingDoubleAt}{71}{78}{CopyingDoubleAt.right_rewrite}
\par\penalty0

The right side's retained tree still requests copying in exactly one search transition at every nonzero original input.
For the proof, the defined maps and incidence equations are substituted.
\lean{Copying/CopyingDoubleAt}{82}{90}{CopyingDoubleAt.right_request}
\par\penalty0

Both sides of the non-distinguished-input diagram take exactly two original machine transitions, irrespective of payload or environment size.
For the proof, the preceding component lemmas are composed or specialized at the stated parameters.
\lean{Copying/CopyingDoubleAt}{94}{102}{CopyingDoubleAt.entry}
\par\penalty0

\subsection{Copying Double At Residual}
The arbitrary-input copying residual is the original copying pair in the retained branch, followed by the independently certified two-copy interchange.

The actual order-exchange filling is supplied to the purely structural assembly; it is used only to reverse the sharing correction, never to assert symmetry of the directed copying relation.
\lean{Copying/CopyingDoubleAtResidual}{19}{27}{CopyingDoubleAt.residual}
\par\penalty0

\subsection{Copying Double Components}
Complete actual occurrences of all components of two literal copies.

Both branches occur before the unchanged outer distributor.
\lean{Copying/CopyingDoubleComponents}{14}{31}{CopyingDouble.Components.pair}
\par\penalty0

The branch map is a complete occurrence, preserving original labels and every ordered source and target attachment.
\lean{Copying/CopyingDoubleComponents}{33}{39}{CopyingDouble.Components.branch}
\par\penalty0

The first map is a complete occurrence, preserving original labels and every ordered source and target attachment.
\lean{Copying/CopyingDoubleComponents}{41}{41}{CopyingDouble.Components.first}
\par\penalty0

The retained map is a complete occurrence, preserving original labels and every ordered source and target attachment.
\lean{Copying/CopyingDoubleComponents}{42}{42}{CopyingDouble.Components.retained}
\par\penalty0

The second active payload has its original nested occurrence.
\lean{Copying/CopyingDoubleComponents}{45}{47}{CopyingDouble.Components.second}
\par\penalty0

The remaining map is a complete occurrence, preserving original labels and every ordered source and target attachment.
\lean{Copying/CopyingDoubleComponents}{49}{52}{CopyingDouble.Components.remaining}
\par\penalty0

The third map is a complete occurrence, preserving original labels and every ordered source and target attachment.
\lean{Copying/CopyingDoubleComponents}{54}{55}{CopyingDouble.Components.third}
\par\penalty0

The tree map is a complete occurrence, preserving original labels and every ordered source and target attachment.
\lean{Copying/CopyingDoubleComponents}{57}{59}{CopyingDouble.Components.tree}
\par\penalty0

The outer Distributor map is a complete occurrence, preserving original labels and every ordered source and target attachment.
\lean{Copying/CopyingDoubleComponents}{61}{62}{CopyingDouble.Components.outerDistributor}
\par\penalty0

The inner Distributor map is a complete occurrence, preserving original labels and every ordered source and target attachment.
\lean{Copying/CopyingDoubleComponents}{64}{65}{CopyingDouble.Components.innerDistributor}
\par\penalty0

The pair input left assertion computes the indicated ordered boundary coordinate in the assembled graph.
For the proof, the defined maps and incidence equations are substituted.
\lean{Copying/CopyingDoubleComponents}{68}{71}{CopyingDouble.Components.pair_input_left}
\par\penalty0

The pair input right assertion computes the indicated ordered boundary coordinate in the assembled graph.
For the proof, the remaining natural-number index bounds follow by linear arithmetic; the defined maps and incidence equations are substituted.
\lean{Copying/CopyingDoubleComponents}{73}{80}{CopyingDouble.Components.pair_input_right}
\par\penalty0

The pair output assertion computes the indicated ordered boundary coordinate in the assembled graph.
For the proof, constructor and boundary cases are separated explicitly; the defined maps and incidence equations are substituted.
\lean{Copying/CopyingDoubleComponents}{82}{87}{CopyingDouble.Components.pair_output}
\par\penalty0

The first input assertion computes the indicated ordered boundary coordinate in the assembled graph.
For the proof, the defined maps and incidence equations are substituted.
\lean{Copying/CopyingDoubleComponents}{89}{91}{CopyingDouble.Components.first_input}
\par\penalty0

The retained input assertion computes the indicated ordered boundary coordinate in the assembled graph.
For the proof, the remaining natural-number index bounds follow by linear arithmetic; the defined maps and incidence equations are substituted.
\lean{Copying/CopyingDoubleComponents}{93}{96}{CopyingDouble.Components.retained_input}
\par\penalty0

The second input assertion computes the indicated ordered boundary coordinate in the assembled graph.
For the proof, the defined maps and incidence equations are substituted.
\lean{Copying/CopyingDoubleComponents}{98}{105}{CopyingDouble.Components.second_input}
\par\penalty0

The tree input assertion computes the indicated ordered boundary coordinate in the assembled graph.
For the proof, the remaining natural-number index bounds follow by linear arithmetic; the defined maps and incidence equations are substituted.
\lean{Copying/CopyingDoubleComponents}{107}{111}{CopyingDouble.Components.tree_input}
\par\penalty0

The tree output assertion computes the indicated ordered boundary coordinate in the assembled graph.
For the proof, the preceding component lemmas are composed or specialized at the stated parameters.
\lean{Copying/CopyingDoubleComponents}{113}{115}{CopyingDouble.Components.tree_output}
\par\penalty0

\subsection{Copying Double Frame}
An explicit common three-payload frame for the two orders of copying. The two active payloads retain inputs zero and one, the residual tree retains all remaining inputs, and each resource has one ordered three-input row slot.

This defines the Resource data for the construction just described. Boundary positions are kept distinct from newly inserted internal vertices.
\lean{Copying/CopyingDoubleFrame}{10}{10}{CopyingDouble.Frame.Resource}
\par\penalty0

This defines the Internal data for the construction just described. Boundary positions are kept distinct from newly inserted internal vertices.
\lean{Copying/CopyingDoubleFrame}{11}{11}{CopyingDouble.Frame.Internal}
\par\penalty0

This defines the Edge data for the construction just described. Boundary positions are kept distinct from newly inserted internal vertices.
\lean{Copying/CopyingDoubleFrame}{12}{12}{CopyingDouble.Frame.Edge}
\par\penalty0

This defines the Vertex data for the construction just described. Boundary positions are kept distinct from newly inserted internal vertices.
\lean{Copying/CopyingDoubleFrame}{13}{13}{CopyingDouble.Frame.Vertex}
\par\penalty0

This defines the payload Input data for the construction just described. Boundary positions are kept distinct from newly inserted internal vertices.
\lean{Copying/CopyingDoubleFrame}{15}{16}{CopyingDouble.Frame.payloadInput}
\par\penalty0

This defines the remaining Input data for the construction just described. Boundary positions are kept distinct from newly inserted internal vertices.
\lean{Copying/CopyingDoubleFrame}{17}{17}{CopyingDouble.Frame.remainingInput}
\par\penalty0

This defines the payload Output data for the construction just described. Boundary positions are kept distinct from newly inserted internal vertices.
\lean{Copying/CopyingDoubleFrame}{18}{18}{CopyingDouble.Frame.payloadOutput}
\par\penalty0

This defines the row Input data for the construction just described. Boundary positions are kept distinct from newly inserted internal vertices.
\lean{Copying/CopyingDoubleFrame}{19}{19}{CopyingDouble.Frame.rowInput}
\par\penalty0

This defines the hole data for the construction just described. Boundary positions are kept distinct from newly inserted internal vertices.
\lean{Copying/CopyingDoubleFrame}{21}{24}{CopyingDouble.Frame.hole}
\par\penalty0

This defines the ordered source list for the construction just described. Boundary positions are kept distinct from newly inserted internal vertices.
\lean{Copying/CopyingDoubleFrame}{25}{28}{CopyingDouble.Frame.source}
\par\penalty0

This defines the ordered target list for the construction just described. Boundary positions are kept distinct from newly inserted internal vertices.
\lean{Copying/CopyingDoubleFrame}{29}{32}{CopyingDouble.Frame.target}
\par\penalty0

This defines the internal wire label for the construction just described. Boundary positions are kept distinct from newly inserted internal vertices.
\lean{Copying/CopyingDoubleFrame}{33}{35}{CopyingDouble.Frame.internalLabel}
\par\penalty0

The triple assertion is the compatibility law for the indicated component of this construction. It relates the actual finite coordinates or certificates, with the hypotheses stated in the linked declaration.
For the proof, constructor and boundary cases are separated explicitly; the remaining natural-number index bounds follow by linear arithmetic.
\lean{Copying/CopyingDoubleFrame}{37}{42}{CopyingDouble.Frame.triple}
\par\penalty0

The payload Input injective assertion states injectivity of the indicated coordinate map: equal image coordinates have equal original coordinates.
For the proof, constructor and boundary cases are separated explicitly; the defined maps and incidence equations are substituted.
\lean{Copying/CopyingDoubleFrame}{44}{48}{CopyingDouble.Frame.payloadInput_injective}
\par\penalty0

The remaining Input injective assertion states injectivity of the indicated coordinate map: equal image coordinates have equal original coordinates.
For the proof, the remaining natural-number index bounds follow by linear arithmetic; the defined maps and incidence equations are substituted.
\lean{Copying/CopyingDoubleFrame}{50}{53}{CopyingDouble.Frame.remainingInput_injective}
\par\penalty0

The payload Output injective assertion states injectivity of the indicated coordinate map: equal image coordinates have equal original coordinates.
For the proof, the preceding component lemmas are composed or specialized at the stated parameters.
\lean{Copying/CopyingDoubleFrame}{55}{57}{CopyingDouble.Frame.payloadOutput_injective}
\par\penalty0

The row Input injective assertion states injectivity of the indicated coordinate map: equal image coordinates have equal original coordinates.
For the proof, the preceding component lemmas are composed or specialized at the stated parameters.
\lean{Copying/CopyingDoubleFrame}{59}{61}{CopyingDouble.Frame.rowInput_injective}
\par\penalty0

The mem remaining equivalence proves both implications of the indicated structural or execution condition.
For the proof, the defined maps and incidence equations are substituted.
\lean{Copying/CopyingDoubleFrame}{63}{65}{CopyingDouble.Frame.mem_remaining}
\par\penalty0

The mem payload equivalence proves both implications of the indicated structural or execution condition.
For the proof, the defined maps and incidence equations are substituted.
\lean{Copying/CopyingDoubleFrame}{66}{68}{CopyingDouble.Frame.mem_payload}
\par\penalty0

The mem row equivalence proves both implications of the indicated structural or execution condition.
For the proof, the defined maps and incidence equations are substituted.
\lean{Copying/CopyingDoubleFrame}{69}{71}{CopyingDouble.Frame.mem_row}
\par\penalty0

The payload ne remaining assertion is the compatibility law for the indicated component of this construction. It relates the actual finite coordinates or certificates, with the hypotheses stated in the linked declaration.
For the proof, constructor and boundary cases are separated explicitly; the remaining natural-number index bounds follow by linear arithmetic; the defined maps and incidence equations are substituted.
\lean{Copying/CopyingDoubleFrame}{73}{78}{CopyingDouble.Frame.payload_ne_remaining}
\par\penalty0

The payload ne row assertion is the compatibility law for the indicated component of this construction. It relates the actual finite coordinates or certificates, with the hypotheses stated in the linked declaration.
For the proof, constructor and boundary cases are separated explicitly; the defined maps and incidence equations are substituted.
\lean{Copying/CopyingDoubleFrame}{80}{82}{CopyingDouble.Frame.payload_ne_row}
\par\penalty0

This defines the finite typed net for the construction just described. Boundary positions are kept distinct from newly inserted internal vertices.
\lean{Copying/CopyingDoubleFrame}{84}{212}{CopyingDouble.Frame.net}
\par\penalty0

The common finite frame in the original depth-indexed copying alphabet.
\lean{Copying/CopyingDoubleFrame}{215}{217}{CopyingDouble.Frame.graph}
\par\penalty0

Every output resource has its original sharing sort.
\lean{Copying/CopyingDoubleFrame}{220}{220}{CopyingDouble.Frame.sort}
\par\penalty0

The sort wire assertion is the compatibility law for the indicated component of this construction. It relates the actual finite coordinates or certificates, with the hypotheses stated in the linked declaration.
For the proof, constructor and boundary cases are separated explicitly; the remaining natural-number index bounds follow by linear arithmetic.
\lean{Copying/CopyingDoubleFrame}{222}{225}{CopyingDouble.Frame.sort_wire}
\par\penalty0

\subsection{Copying Double Interchange}
Two original Figure 17 copies commute up to the independently proved sharing-tree template. Both full canonical distributors remain in each graph, all three payload copies are retained, and every external input is fixed.

Copying original inputs zero and one in either order produces a genuine sharing-template filling witness. This theorem includes arbitrary remaining inputs and arbitrary ordered term/atom environments, including empty forests. It asserts no isomorphism between different sharing-tree presentations.
\lean{Copying/CopyingDoubleInterchange}{12}{22}{CopyingDouble.interchange}
\par\penalty0

\subsection{Copying Double Normal}
A common normal presentation with three identical payloads, the full remaining contraction tree, and actual two-layer canonical resource rows.

This defines the row data for the construction just described. Boundary positions are kept distinct from newly inserted internal vertices.
\lean{Copying/CopyingDoubleNormal}{13}{17}{CopyingDouble.Frame.row}
\par\penalty0

This defines the typed family of inserted components for the construction just described. Boundary positions are kept distinct from newly inserted internal vertices.
\lean{Copying/CopyingDoubleNormal}{19}{24}{CopyingDouble.Frame.family}
\par\penalty0

This is a finite expansion of an explicit three-payload frame.
\lean{Copying/CopyingDoubleNormal}{27}{30}{CopyingDouble.Frame.normal}
\par\penalty0

Interchanging two uses changes only independently proved sharing-tree components; every payload is literally retained in the same ordered slot.
\lean{Copying/CopyingDoubleNormal}{34}{44}{CopyingDouble.Frame.witness}
\par\penalty0

\subsection{Copying Double Resources}
Complete six-edge resource rows in the two-copy graph, with original payload-output attachments and no canonical-tree simplification.

This defines the atom Position data for the construction just described. Boundary positions are kept distinct from newly inserted internal vertices.
\lean{Copying/CopyingDoubleResources}{11}{15}{CopyingDouble.Components.atomPosition}
\par\penalty0

The resource position retains its original term/atom sort.
\lean{Copying/CopyingDoubleResources}{18}{67}{CopyingDouble.Components.resource}
\par\penalty0

The resource input assertion computes the indicated ordered boundary coordinate in the assembled graph.
For the proof, constructor and boundary cases are separated explicitly; the remaining natural-number index bounds follow by linear arithmetic; the defined maps and incidence equations are substituted.
\lean{Copying/CopyingDoubleResources}{69}{91}{CopyingDouble.Components.resource_input}
\par\penalty0

The resource output assertion computes the indicated ordered boundary coordinate in the assembled graph.
For the proof, constructor and boundary cases are separated explicitly; the remaining natural-number index bounds follow by linear arithmetic; the defined maps and incidence equations are substituted.
\lean{Copying/CopyingDoubleResources}{93}{102}{CopyingDouble.Components.resource_output}
\par\penalty0

Distinct resource rows retain disjoint original edge carriers.
For the proof, constructor and boundary cases are separated explicitly; the remaining natural-number index bounds follow by linear arithmetic; the defined maps and incidence equations are substituted.
\lean{Copying/CopyingDoubleResources}{105}{120}{CopyingDouble.Components.resource_disjoint}
\par\penalty0

The original environment distributor is exhausted by its actual rows.
For the proof, constructor and boundary cases are separated explicitly; the defined maps and incidence equations are substituted.
\lean{Copying/CopyingDoubleResources}{123}{139}{CopyingDouble.Components.resource_cover}
\par\penalty0

The inner Row map is a complete occurrence, preserving original labels and every ordered source and target attachment.
\lean{Copying/CopyingDoubleResources}{141}{143}{CopyingDouble.Components.innerRow}
\par\penalty0

The outer Row map is a complete occurrence, preserving original labels and every ordered source and target attachment.
\lean{Copying/CopyingDoubleResources}{144}{146}{CopyingDouble.Components.outerRow}
\par\penalty0

The row link assertion is the compatibility law for the indicated component of this construction. It relates the actual finite coordinates or certificates, with the hypotheses stated in the linked declaration.
For the proof, the defined maps and incidence equations are substituted.
\lean{Copying/CopyingDoubleResources}{148}{154}{CopyingDouble.Components.row_link}
\par\penalty0

The row disjoint assertion excludes a shared edge or coordinate between the indicated independently inserted components.
For the proof, constructor and boundary cases are separated explicitly.
\lean{Copying/CopyingDoubleResources}{156}{158}{CopyingDouble.Components.row_disjoint}
\par\penalty0

Both nested original distributors form one actual connected six-edge row.
\lean{Copying/CopyingDoubleResources}{161}{164}{CopyingDouble.Components.row}
\par\penalty0

The outer input assertion computes the indicated ordered boundary coordinate in the assembled graph.
For the proof, the defined maps and incidence equations are substituted.
\lean{Copying/CopyingDoubleResources}{167}{172}{CopyingDouble.Components.outer_input}
\par\penalty0

The inner input active assertion computes the indicated ordered boundary coordinate in the assembled graph.
For the proof, the defined maps and incidence equations are substituted.
\lean{Copying/CopyingDoubleResources}{174}{181}{CopyingDouble.Components.inner_input_active}
\par\penalty0

The inner input retained assertion computes the indicated ordered boundary coordinate in the assembled graph.
For the proof, the defined maps and incidence equations are substituted.
\lean{Copying/CopyingDoubleResources}{183}{192}{CopyingDouble.Components.inner_input_retained}
\par\penalty0

All three row inputs remain attached to the original corresponding payload output; these are exact coordinates, not just unlabelled paths.
For the proof, constructor and boundary cases are separated explicitly; the defined maps and incidence equations are substituted.
\lean{Copying/CopyingDoubleResources}{196}{202}{CopyingDouble.Components.row_input}
\par\penalty0

The row output assertion computes the indicated ordered boundary coordinate in the assembled graph.
For the proof, the defined maps and incidence equations are substituted.
\lean{Copying/CopyingDoubleResources}{204}{207}{CopyingDouble.Components.row_output}
\par\penalty0

The rows disjoint assertion excludes a shared edge or coordinate between the indicated independently inserted components.
For the proof, constructor and boundary cases are separated explicitly.
\lean{Copying/CopyingDoubleResources}{210}{220}{CopyingDouble.Components.rows_disjoint}
\par\penalty0

\subsection{Copying Double Rewrite}
The second original copying step occurs in the retained branch of the literal first reduct. The outer active copy and canonical distributors remain.

This defines the retained Context data for the construction just described. Boundary positions are kept distinct from newly inserted internal vertices.
\lean{Copying/CopyingDoubleRewrite}{13}{14}{CopyingDouble.retainedContext}
\par\penalty0

This defines the retained Hole data for the construction just described. Boundary positions are kept distinct from newly inserted internal vertices.
\lean{Copying/CopyingDoubleRewrite}{16}{17}{CopyingDouble.retainedHole}
\par\penalty0

The retained Iso map is an ordered recursive graph isomorphism between the stated original constructions.
\lean{Copying/CopyingDoubleRewrite}{19}{31}{CopyingDouble.retainedIso}
\par\penalty0

A real retained-branch rewrite keeps the active copy and preserves the original external positions: retained input i is external input i+1.
For the proof, constructor and boundary cases are separated explicitly; the remaining natural-number index bounds follow by linear arithmetic; the defined maps and incidence equations are substituted.
\lean{Copying/CopyingDoubleRewrite}{35}{54}{CopyingDouble.copiesWith_rewrite}
\par\penalty0

The surrounding original resource distributor is also retained by the compute step, using the machine's actual simultaneous-expansion closure.
For the proof, constructor and boundary cases are separated explicitly; the remaining natural-number index bounds follow by linear arithmetic; the defined maps and incidence equations are substituted.
\lean{Copying/CopyingDoubleRewrite}{58}{72}{CopyingDouble.afterWith_rewrite}
\par\penalty0

Copying original input one of the displayed reduct is exactly a second original Figure 17 rewrite, with no algebraic contraction step inserted.
For the proof, the defined maps and incidence equations are substituted.
\lean{Copying/CopyingDoubleRewrite}{76}{80}{CopyingDouble.second_rewrite}
\par\penalty0

\subsection{Copying Double Rows}
The exact resource row produced by two nested literal copying steps. It contains both complete D1,2 trees, with their weakenings and ordered incidences. Exchanging the first two payloads changes only the ordered row inputs.

This defines the common graph frame for the construction just described. Boundary positions are kept distinct from newly inserted internal vertices.
\lean{Copying/CopyingDoubleRows}{12}{15}{CopyingDouble.Row.frame}
\par\penalty0

This defines the typed family of inserted components for the construction just described. Boundary positions are kept distinct from newly inserted internal vertices.
\lean{Copying/CopyingDoubleRows}{17}{20}{CopyingDouble.Row.family}
\par\penalty0

The inner row merges payloads one and two; its result is the second input of the outer row, whose first input comes from payload zero.
\lean{Copying/CopyingDoubleRows}{24}{25}{CopyingDouble.Row.graph}
\par\penalty0

The sharing assertion is the compatibility law for the indicated component of this construction. It relates the actual finite coordinates or certificates, with the hypotheses stated in the linked declaration.
For the proof, constructor and boundary cases are separated explicitly.
\lean{Copying/CopyingDoubleRows}{27}{32}{CopyingDouble.Row.sharing}
\par\penalty0

A permutation changes actual row incidences and retains tree acyclicity.
For the proof, the preceding component lemmas are composed or specialized at the stated parameters.
\lean{Copying/CopyingDoubleRows}{35}{41}{CopyingDouble.Row.permuted_sharing}
\par\penalty0

This defines the exchange data for the construction just described. Boundary positions are kept distinct from newly inserted internal vertices.
\lean{Copying/CopyingDoubleRows}{43}{44}{CopyingDouble.Row.exchange}
\par\penalty0

Both nested rows are a literal proved sharing-tree relation, not an assumption that resource distributors commute.
For the proof, the preceding component lemmas are composed or specialized at the stated parameters.
\lean{Copying/CopyingDoubleRows}{48}{50}{CopyingDouble.Row.relation}
\par\penalty0

The witness data give an actual recursive paired filling, with its common context and endpoint representation maps.
\lean{Copying/CopyingDoubleRows}{52}{57}{CopyingDouble.Row.witness}
\par\penalty0

Coordinates of the exact nested two-row frame.
\lean{Copying/CopyingDoubleRows}{60}{65}{CopyingDouble.Row.host}
\par\penalty0

The components map is a complete occurrence, preserving original labels and every ordered source and target attachment.
\lean{Copying/CopyingDoubleRows}{67}{71}{CopyingDouble.Row.components}
\par\penalty0

Actual nested D1,2 occurrences assemble to the full three-input row. No vertex disjointness is postulated: complete incidence establishes it.
\lean{Copying/CopyingDoubleRows}{76}{103}{CopyingDouble.Row.occurrence}
\par\penalty0

\subsection{Copying Double Search}
Exact search requests at the non-distinguished inputs of the original copying reduct. Search moves through genuine context attachment maps only.

The step reboundary assertion supplies the indicated actual machine transition or its transport through the common construction.
For the proof, constructor and boundary cases are separated explicitly; the defined maps and incidence equations are substituted.
\lean{Copying/CopyingDoubleSearch}{10}{22}{Search.step_reboundary}
\par\penalty0

The copies With request assertion is the compatibility law for the indicated component of this construction. It relates the actual finite coordinates or certificates, with the hypotheses stated in the linked declaration.
For the proof, constructor and boundary cases are separated explicitly; the remaining natural-number index bounds follow by linear arithmetic; the defined maps and incidence equations are substituted.
\lean{Copying/CopyingDoubleSearch}{27}{47}{CopyingDouble.copiesWith_request}
\par\penalty0

The after With request assertion is the compatibility law for the indicated component of this construction. It relates the actual finite coordinates or certificates, with the hypotheses stated in the linked declaration.
For the proof, constructor and boundary cases are separated explicitly; the remaining natural-number index bounds follow by linear arithmetic; the defined maps and incidence equations are substituted.
\lean{Copying/CopyingDoubleSearch}{49}{62}{CopyingDouble.afterWith_request}
\par\penalty0

Original input one requests and executes exactly one further original copy.
For the proof, the defined maps and incidence equations are substituted.
\lean{Copying/CopyingDoubleSearch}{65}{71}{CopyingDouble.second_entry}
\par\penalty0

\subsection{Copying Double Source}
Exact normalization of the original two-copy target. A complete occurrence of the common frame exhausts the actual original edges, so its isomorphism is proved from incidence and finite carrier coverage.

The payload map is a complete occurrence, preserving original labels and every ordered source and target attachment.
\lean{Copying/CopyingDoubleSource}{13}{15}{CopyingDouble.Source.payload}
\par\penalty0

The each map is a complete occurrence, preserving original labels and every ordered source and target attachment.
\lean{Copying/CopyingDoubleSource}{17}{20}{CopyingDouble.Source.each}
\par\penalty0

This defines the old-frame port map for the construction just described. Boundary positions are kept distinct from newly inserted internal vertices.
\lean{Copying/CopyingDoubleSource}{22}{26}{CopyingDouble.Source.host}
\par\penalty0

The input assertion computes the indicated ordered boundary coordinate in the assembled graph.
For the proof, constructor and boundary cases are separated explicitly; the defined maps and incidence equations are substituted.
\lean{Copying/CopyingDoubleSource}{28}{44}{CopyingDouble.Source.input}
\par\penalty0

The output assertion computes the indicated ordered boundary coordinate in the assembled graph.
For the proof, constructor and boundary cases are separated explicitly; the defined maps and incidence equations are substituted.
\lean{Copying/CopyingDoubleSource}{46}{61}{CopyingDouble.Source.output}
\par\penalty0

The payloads disjoint assertion excludes a shared edge or coordinate between the indicated independently inserted components.
For the proof, constructor and boundary cases are separated explicitly; the defined maps and incidence equations are substituted.
\lean{Copying/CopyingDoubleSource}{64}{71}{CopyingDouble.Source.payloads_disjoint}
\par\penalty0

The payload tree disjoint assertion excludes a shared edge or coordinate between the indicated independently inserted components.
For the proof, constructor and boundary cases are separated explicitly; the defined maps and incidence equations are substituted.
\lean{Copying/CopyingDoubleSource}{73}{78}{CopyingDouble.Source.payload_tree_disjoint}
\par\penalty0

The payload row disjoint assertion excludes a shared edge or coordinate between the indicated independently inserted components.
For the proof, constructor and boundary cases are separated explicitly; the defined maps and incidence equations are substituted.
\lean{Copying/CopyingDoubleSource}{80}{85}{CopyingDouble.Source.payload_row_disjoint}
\par\penalty0

The tree row disjoint assertion excludes a shared edge or coordinate between the indicated independently inserted components.
For the proof, constructor and boundary cases are separated explicitly.
\lean{Copying/CopyingDoubleSource}{87}{90}{CopyingDouble.Source.tree_row_disjoint}
\par\penalty0

The disjoint assertion excludes a shared edge or coordinate between the indicated independently inserted components.
For the proof, constructor and boundary cases are separated explicitly.
\lean{Copying/CopyingDoubleSource}{92}{106}{CopyingDouble.Source.disjoint}
\par\penalty0

The common normal graph is a complete-label occurrence in the actual original target, with the original external boundary fixed pointwise.
\lean{Copying/CopyingDoubleSource}{110}{112}{CopyingDouble.Source.occurrence}
\par\penalty0

Every original edge belongs to exactly the retained payload, remaining tree, or nested-resource component described by the common frame.
For the proof, constructor and boundary cases are separated explicitly; the defined maps and incidence equations are substituted.
\lean{Copying/CopyingDoubleSource}{117}{144}{CopyingDouble.Source.cover}
\par\penalty0

Exact original-graph normalization. All canonical rows, three payload copies, and both inserted weakenings of the remaining tree are preserved.
\lean{Copying/CopyingDoubleSource}{148}{150}{CopyingDouble.Source.iso}
\par\penalty0

\subsection{Copying Double Symmetry}
Exchanging the two active payload slots in the explicit common frame. The actual row input permutation and outer input compensation are both retained.

Original input reindexing gives a complete ordered occurrence.
\lean{Copying/CopyingDoubleSymmetry}{10}{21}{CopyingDouble.Symmetry.forward}
\par\penalty0

Reverse the same actual ordered input action.
\lean{Copying/CopyingDoubleSymmetry}{24}{41}{CopyingDouble.Symmetry.backward}
\par\penalty0

This defines the swap data for the construction just described. Boundary positions are kept distinct from newly inserted internal vertices.
\lean{Copying/CopyingDoubleSymmetry}{43}{43}{CopyingDouble.Symmetry.swap}
\par\penalty0

The swap swap assertion is the compatibility law for the indicated component of this construction. It relates the actual finite coordinates or certificates, with the hypotheses stated in the linked declaration.
For the proof, the preceding component lemmas are composed or specialized at the stated parameters.
\lean{Copying/CopyingDoubleSymmetry}{45}{46}{CopyingDouble.Symmetry.swap_swap}
\par\penalty0

The swap injective assertion states injectivity of the indicated coordinate map: equal image coordinates have equal original coordinates.
For the proof, the preceding component lemmas are composed or specialized at the stated parameters.
\lean{Copying/CopyingDoubleSymmetry}{48}{53}{CopyingDouble.Symmetry.swap_injective}
\par\penalty0

The row Back map is a complete occurrence, preserving original labels and every ordered source and target attachment.
\lean{Copying/CopyingDoubleSymmetry}{55}{59}{CopyingDouble.Symmetry.rowBack}
\par\penalty0

The row Back input assertion computes the indicated ordered boundary coordinate in the assembled graph.
For the proof, constructor and boundary cases are separated explicitly; the defined maps and incidence equations are substituted.
\lean{Copying/CopyingDoubleSymmetry}{61}{63}{CopyingDouble.Symmetry.rowBack_input}
\par\penalty0

The row Back output assertion computes the indicated ordered boundary coordinate in the assembled graph.
For the proof, the defined maps and incidence equations are substituted.
\lean{Copying/CopyingDoubleSymmetry}{65}{68}{CopyingDouble.Symmetry.rowBack_output}
\par\penalty0

Exchange only the two active payload slots; retain the full residual tree and every original resource row slot.
\lean{Copying/CopyingDoubleSymmetry}{72}{74}{CopyingDouble.Symmetry.slot}
\par\penalty0

The slot slot assertion is the compatibility law for the indicated component of this construction. It relates the actual finite coordinates or certificates, with the hypotheses stated in the linked declaration.
For the proof, constructor and boundary cases are separated explicitly; the defined maps and incidence equations are substituted.
\lean{Copying/CopyingDoubleSymmetry}{76}{79}{CopyingDouble.Symmetry.slot_slot}
\par\penalty0

The slot injective assertion states injectivity of the indicated coordinate map: equal image coordinates have equal original coordinates.
For the proof, the preceding component lemmas are composed or specialized at the stated parameters.
\lean{Copying/CopyingDoubleSymmetry}{81}{84}{CopyingDouble.Symmetry.slot_injective}
\par\penalty0

The component map is a complete occurrence, preserving original labels and every ordered source and target attachment.
\lean{Copying/CopyingDoubleSymmetry}{86}{89}{CopyingDouble.Symmetry.component}
\par\penalty0

The each map is a complete occurrence, preserving original labels and every ordered source and target attachment.
\lean{Copying/CopyingDoubleSymmetry}{91}{95}{CopyingDouble.Symmetry.each}
\par\penalty0

This defines the old-frame port map for the construction just described. Boundary positions are kept distinct from newly inserted internal vertices.
\lean{Copying/CopyingDoubleSymmetry}{97}{102}{CopyingDouble.Symmetry.host}
\par\penalty0

The input assertion computes the indicated ordered boundary coordinate in the assembled graph.
For the proof, constructor and boundary cases are separated explicitly; the remaining natural-number index bounds follow by linear arithmetic; the defined maps and incidence equations are substituted.
\lean{Copying/CopyingDoubleSymmetry}{105}{131}{CopyingDouble.Symmetry.input}
\par\penalty0

The output assertion computes the indicated ordered boundary coordinate in the assembled graph.
For the proof, constructor and boundary cases are separated explicitly; the defined maps and incidence equations are substituted.
\lean{Copying/CopyingDoubleSymmetry}{133}{146}{CopyingDouble.Symmetry.output}
\par\penalty0

The disjoint assertion excludes a shared edge or coordinate between the indicated independently inserted components.
For the proof, the preceding component lemmas are composed or specialized at the stated parameters.
\lean{Copying/CopyingDoubleSymmetry}{148}{152}{CopyingDouble.Symmetry.disjoint}
\par\penalty0

Actual complete occurrence into the input-compensated opposite frame.
\lean{Copying/CopyingDoubleSymmetry}{155}{157}{CopyingDouble.Symmetry.occurrence}
\par\penalty0

The cover assertion is the compatibility law for the indicated component of this construction. It relates the actual finite coordinates or certificates, with the hypotheses stated in the linked declaration.
For the proof, constructor and boundary cases are separated explicitly; the defined maps and incidence equations are substituted.
\lean{Copying/CopyingDoubleSymmetry}{160}{170}{CopyingDouble.Symmetry.cover}
\par\penalty0

The fixed-boundary symmetry exchanges the first two identical payloads, retains the entire remaining tree, and reverses the two resource-row inputs.
\lean{Copying/CopyingDoubleSymmetry}{174}{176}{CopyingDouble.Symmetry.iso}
\par\penalty0

\subsection{Copying Double Witness}
Actual residual copying certificates after entering another original tree input. Unrelated branches and the canonical resource distributors remain in the shared context.

Assemble any actual paired retained branch beneath the original active payload and the two-use term/atom distributors.
\lean{Copying/CopyingDoubleWitness}{13}{26}{CopyingDouble.afterWithWitness}
\par\penalty0

After copying input zero, copying input one leaves one genuine copying pair in the retained branch; the active first payload is part of the context.
\lean{Copying/CopyingDoubleWitness}{30}{35}{CopyingDouble.secondWitness}
\par\penalty0

The opposite order has the same literal residual copying certificate, transported through the actual input-swap context.
\lean{Copying/CopyingDoubleWitness}{39}{48}{CopyingDouble.oppositeWitness}
\par\penalty0

\subsection{Copying Filling Reverse}
Reverse a genuine recursive filling only when its underlying template relation has a proved symmetry. This applies to the sharing correction, not to the directed copying pre-template itself.

The reverse assertion is the compatibility law for the indicated component of this construction. It relates the actual finite coordinates or certificates, with the hypotheses stated in the linked declaration.
For the proof, induction follows the original derivation or finite structure; constructor and boundary cases are separated explicitly.
\lean{Copying/CopyingFillingReverse}{11}{31}{DeepFilling.reverse}
\par\penalty0

The reverse data give an actual recursive paired filling, with its common context and endpoint representation maps.
\lean{Copying/CopyingFillingReverse}{33}{36}{FillingWitness.reverse}
\par\penalty0

The reverse Sharing data give an actual recursive paired filling, with its common context and endpoint representation maps.
\lean{Copying/CopyingFillingReverse}{38}{41}{FillingWitness.reverseSharing}
\par\penalty0

\subsection{Copying Input Context}
Input permutation is an actual one-hole graph context. Consequently it transports original rewrites with their exact input positions and rule premises.

This defines the common graph frame for the construction just described. Boundary positions are kept distinct from newly inserted internal vertices.
\lean{Copying/CopyingInputContext}{10}{11}{CopyingInputPermutation.frame}
\par\penalty0

This defines the frame Edge data for the construction just described. Boundary positions are kept distinct from newly inserted internal vertices.
\lean{Copying/CopyingInputContext}{13}{14}{CopyingInputPermutation.frameEdge}
\par\penalty0

The plugged payload uses precisely the original interfaces of the hole.
\lean{Copying/CopyingInputContext}{17}{19}{CopyingInputPermutation.inFrame}
\par\penalty0

Delete only the empty corolla carrier introduced by the permutation context.
\lean{Copying/CopyingInputContext}{22}{53}{CopyingInputPermutation.frameIso}
\par\penalty0

Reindexing input ports preserves each actual original rewrite and moves the active and resumed positions by the same typed permutation.
For the proof, constructor and boundary cases are separated explicitly.
\lean{Copying/CopyingInputContext}{57}{69}{CopyingInputPermutation.rewrite}
\par\penalty0

\subsection{Copying Input Isomorphism}
Fixed-interface compensation for the actual input-permutation action.

Cancel a recorded input action while preserving every ordered output.
\lean{Copying/CopyingInputIsomorphism}{10}{30}{CopyingInputPermutation.compensate}
\par\penalty0

The roundtrip map is an ordered recursive graph isomorphism between the stated original constructions.
\lean{Copying/CopyingInputIsomorphism}{32}{34}{CopyingInputPermutation.roundtrip}
\par\penalty0

The congr map is an ordered recursive graph isomorphism between the stated original constructions.
\lean{Copying/CopyingInputIsomorphism}{36}{50}{CopyingInputPermutation.congr}
\par\penalty0

At every depth the reconstructed tree retains exactly the permuted original incidences and full primitive labels.
\lean{Copying/CopyingInputIsomorphism}{54}{73}{CopyingInputPermutation.reorderedIso}
\par\penalty0

Reordering a complete contraction tree and compensating its external input action gives the same original copying source, including the original payload.
\lean{Copying/CopyingInputIsomorphism}{77}{85}{CopyingInputPermutation.beforeIso}
\par\penalty0

\subsection{Copying Input Permutation}
Actual interface permutation: move boundary occurrences, preserving all carriers, edge labels and incidence order. This is not a new term translation.

This defines the port map for the construction just described. Boundary positions are kept distinct from newly inserted internal vertices.
\lean{Copying/CopyingInputPermutation}{10}{20}{CopyingInputPermutation.ports}
\par\penalty0

The input iff equivalence proves both implications of the indicated structural or execution condition.
For the proof, constructor and boundary cases are separated explicitly; the defined maps and incidence equations are substituted.
\lean{Copying/CopyingInputPermutation}{22}{22}{CopyingInputPermutation.input_iff}
\par\penalty0

The output iff equivalence proves both implications of the indicated structural or execution condition.
For the proof, constructor and boundary cases are separated explicitly; the defined maps and incidence equations are substituted.
\lean{Copying/CopyingInputPermutation}{23}{23}{CopyingInputPermutation.output_iff}
\par\penalty0

The label assertion verifies the indicated label compatibility or label exclusion at the original coordinate.
For the proof, constructor and boundary cases are separated explicitly; the defined maps and incidence equations are substituted.
\lean{Copying/CopyingInputPermutation}{25}{29}{CopyingInputPermutation.label}
\par\penalty0

Forward-position convention: old input i moves to input p(i). The source incidences retain every original source-port order.
\lean{Copying/CopyingInputPermutation}{33}{87}{CopyingInputPermutation.net}
\par\penalty0

The ordered incidence isomorphism records, rather than forgets, the external action.
\lean{Copying/CopyingInputPermutation}{90}{98}{CopyingInputPermutation.iso}
\par\penalty0

The iso output assertion computes the indicated ordered boundary coordinate in the assembled graph.
Both sides are definitionally equal.
\lean{Copying/CopyingInputPermutation}{100}{100}{CopyingInputPermutation.iso_output}
\par\penalty0

The iso input assertion computes the indicated ordered boundary coordinate in the assembled graph.
Both sides are definitionally equal.
\lean{Copying/CopyingInputPermutation}{101}{101}{CopyingInputPermutation.iso_input}
\par\penalty0

The all Labels iff equivalence proves both implications of the indicated structural or execution condition.
For the proof, the defined maps and incidence equations are substituted.
\lean{Copying/CopyingInputPermutation}{103}{103}{CopyingInputPermutation.allLabels_iff}
\par\penalty0

The binding Free iff equivalence proves both implications of the indicated structural or execution condition.
For the proof, the preceding component lemmas are composed or specialized at the stated parameters.
\lean{Copying/CopyingInputPermutation}{105}{105}{CopyingInputPermutation.bindingFree_iff}
\par\penalty0

\subsection{Copying Input Tree}
Reindex the full ordered input interface of a genuine contraction tree. The selected port becomes zero; all other original edges and ports remain.

This defines the exchange data for the construction just described. Boundary positions are kept distinct from newly inserted internal vertices.
\lean{Copying/CopyingInputTree}{13}{14}{CopyingInputPermutation.exchange}
\par\penalty0

The exchange involutive assertion is the compatibility law for the indicated component of this construction. It relates the actual finite coordinates or certificates, with the hypotheses stated in the linked declaration.
For the proof, constructor and boundary cases are separated explicitly; the defined maps and incidence equations are substituted.
\lean{Copying/CopyingInputTree}{16}{24}{CopyingInputPermutation.exchange_involutive}
\par\penalty0

This defines the swap data for the construction just described. Boundary positions are kept distinct from newly inserted internal vertices.
\lean{Copying/CopyingInputTree}{26}{31}{CopyingInputPermutation.swap}
\par\penalty0

The reordered graph assertion is the compatibility law for the indicated component of this construction. It relates the actual finite coordinates or certificates, with the hypotheses stated in the linked declaration.
For the proof, induction follows the original derivation or finite structure; constructor and boundary cases are separated explicitly; the defined maps and incidence equations are substituted.
\lean{Copying/CopyingInputTree}{33}{50}{CopyingInputPermutation.reordered_graph}
\par\penalty0

This defines the reordered data for the construction just described. Boundary positions are kept distinct from newly inserted internal vertices.
\lean{Copying/CopyingInputTree}{52}{54}{CopyingInputPermutation.reordered}
\par\penalty0

The reordered label assertion verifies the indicated label compatibility or label exclusion at the original coordinate.
For the proof, the preceding component lemmas are composed or specialized at the stated parameters.
\lean{Copying/CopyingInputTree}{56}{60}{CopyingInputPermutation.reordered_label}
\par\penalty0

The reindexed tree has the same edge carrier. Its occurrence maps the new input positions back to the original ordered inputs using the inverse.
\lean{Copying/CopyingInputTree}{64}{87}{CopyingInputPermutation.occurrence}
\par\penalty0

The occurrence input assertion computes the indicated ordered boundary coordinate in the assembled graph.
Both sides are definitionally equal.
\lean{Copying/CopyingInputTree}{89}{92}{CopyingInputPermutation.occurrence_input}
\par\penalty0

The occurrence output assertion computes the indicated ordered boundary coordinate in the assembled graph.
For the proof, the defined maps and incidence equations are substituted.
\lean{Copying/CopyingInputTree}{94}{96}{CopyingInputPermutation.occurrence_output}
\par\penalty0

\subsection{Copying Input Witness}
The actual input-permutation graph context transports recursive paired fillings without changing the underlying template relation.

Relabelled external input positions are part of the common graph context; they require no extra constructor or closure assumption on the pre-template.
\lean{Copying/CopyingInputWitness}{13}{22}{FillingWitness.inputPermutation}
\par\penalty0

An arbitrary-input copying step itself has an independently constructed recursive copying-template certificate. The original relation still uses input zero; the actual common context supplies precisely the compensated permutation.
\lean{Copying/CopyingInputWitness}{32}{40}{CopyingSelectedInput.witness}
\par\penalty0

\subsection{Copying Residual Tree}
The residual sharing component of the original copying rule is itself a finite contraction tree. Its selected input is fed by the actual weakening.

The residual Frame acyclic assertion is the compatibility law for the indicated component of this construction. It relates the actual finite coordinates or certificates, with the hypotheses stated in the linked declaration.
For the proof, constructor and boundary cases are separated explicitly.
\lean{Copying/CopyingResidualTree}{11}{15}{CopyRule.residualFrame_acyclic}
\par\penalty0

Feeding the selected input with weakening leaves every remaining input in its original order and preserves genuine tree acyclicity.
\lean{Copying/CopyingResidualTree}{19}{38}{CopyRule.residualTree}
\par\penalty0

The residual Tree at Depth assertion is the compatibility law for the indicated component of this construction. It relates the actual finite coordinates or certificates, with the hypotheses stated in the linked declaration.
Both sides are definitionally equal.
\lean{Copying/CopyingResidualTree}{40}{40}{CopyRule.residualTree_atDepth}
\par\penalty0

\subsection{Copying Selected Input}
Every original input of a copying tree admits an actual Figure 17 rewrite. The finite target retains the complete reordered original tree and restores the external interface; no selected-spine or unused-branch assumption is needed.

This defines the permutation data for the construction just described. Boundary positions are kept distinct from newly inserted internal vertices.
\lean{Copying/CopyingSelectedInput}{16}{17}{CopyingSelectedInput.permutation}
\par\penalty0

Literal Figure 17 reduct for the selected original input, with the external input reindexing cancelled after the rewrite.
\lean{Copying/CopyingSelectedInput}{21}{24}{CopyingSelectedInput.after}
\par\penalty0

The only rewrite used is the original local copying rule, transported through a proved actual input-permutation context and a recursive isomorphism.
For the proof, the defined maps and incidence equations are substituted.
\lean{Copying/CopyingSelectedInput}{28}{38}{CopyingSelectedInput.rewrite}
\par\penalty0

All original inputs have the same exact request/compute prefix: one search transition followed by one genuine copying transition.
For the proof, the preceding component lemmas are composed or specialized at the stated parameters.
\lean{Copying/CopyingSelectedInput}{42}{46}{CopyingSelectedInput.entry}
\par\penalty0

\subsection{Copying Template}
The copying pre-template is the token-erasure of the original Figure 17 rule. The entire ordered input interface and the residual contraction tree are retained, including the selected input's new weakening.

Figure 17 without its token: an arbitrary term contraction tree and the original copyable payload, with the two actual output distributors.
\lean{Copying/CopyingTemplate}{14}{19}{CopyingTemplate.Literal}
\par\penalty0

Each side may use its own graph carrier while preserving all ordered external ports and every recursive box body.
\lean{Copying/CopyingTemplate}{23}{26}{CopyingTemplate.Pair}
\par\penalty0

The focus Free assertion excludes focus labels throughout the indicated graph, including recursively inside boxes.
For the proof, constructor and boundary cases are separated explicitly.
\lean{Copying/CopyingTemplate}{28}{34}{CopyingTemplate.Literal.focusFree}
\par\penalty0

The focus Free assertion excludes focus labels throughout the indicated graph, including recursively inside boxes.
For the proof, constructor and boundary cases are separated explicitly.
\lean{Copying/CopyingTemplate}{36}{40}{CopyingTemplate.Pair.focusFree}
\par\penalty0

This pre-template has no assumed simulation, robustness, or contextual equivalence; its relation is precisely the original copying graphs.
\lean{Copying/CopyingTemplate}{44}{46}{CopyingTemplate.template}
\par\penalty0

The literal assertion excludes focus labels throughout the indicated graph, including recursively inside boxes.
For the proof, the preceding component lemmas are composed or specialized at the stated parameters.
\lean{Copying/CopyingTemplate}{48}{52}{CopyingTemplate.literal}
\par\penalty0

The token-erased rule retains every input, not only the distinguished input at which its compute step is specified.
For the proof, constructor and boundary cases are separated explicitly; the defined maps and incidence equations are substituted.
\lean{Copying/CopyingTemplate}{56}{59}{CopyingTemplate.Literal.inputs}
\par\penalty0

The inputs assertion computes the indicated ordered boundary coordinate in the assembled graph.
For the proof, constructor and boundary cases are separated explicitly.
\lean{Copying/CopyingTemplate}{61}{64}{CopyingTemplate.Pair.inputs}
\par\penalty0

\subsection{Copying Template Entry}
The original copying entry transitions, with all tree inputs retained. The distinguished first input consumes the template in exactly two machine transitions. Other inputs request copying too, but their residual comparison is a separate associativity/commutativity obligation.

Every input of the original tree requests copying in one search step; this includes the inputs that are not distinguished by the displayed rule.
For the proof, the defined maps and incidence equations are substituted.
\lean{Copying/CopyingTemplateEntry}{14}{19}{CopyingTemplate.requests_copy}
\par\penalty0

Independent recursive graph representations preserve the search request at every original external input coordinate.
For the proof, constructor and boundary cases are separated explicitly.
\lean{Copying/CopyingTemplateEntry}{23}{29}{CopyingTemplate.Pair.search_request}
\par\penalty0

The distinguished first input executes exactly the original local copy rule, returning to search at that same input on the actual displayed reduct.
For the proof, constructor and boundary cases are separated explicitly; the defined maps and incidence equations are substituted.
\lean{Copying/CopyingTemplateEntry}{33}{43}{CopyingTemplate.Pair.selected_entry_prefix}
\par\penalty0

At the first input the original entered hole is consumed into the common context, with every other shallow or deep replacement retained. The left endpoint takes two transitions and the right endpoint takes none.
For the proof, the preceding component lemmas are composed or specialized at the stated parameters.
\lean{Copying/CopyingTemplateEntry}{48}{60}{CopyingTemplate.consume_selected}
\par\penalty0

Unrestricted contexts discharge the selected-input consumption without an additional context admissibility premise.
For the proof, the preceding component lemmas are composed or specialized at the stated parameters.
\lean{Copying/CopyingTemplateEntry}{64}{72}{CopyingTemplate.consume_selected_allContexts}
\par\penalty0

The same first-input residual is binding-free by the proved sourced-hole plugging law and the actual closed right component.
For the proof, the preceding component lemmas are composed or specialized at the stated parameters.
\lean{Copying/CopyingTemplateEntry}{76}{88}{CopyingTemplate.consume_selected_bindingFree}
\par\penalty0

\subsection{Copying Template Input Assembly}
Lift the two actual non-distinguished-input prefixes and their original copying/sharing residual into an arbitrary original recursive specimen.

The actual one-search/one-copy prefix is preserved by simultaneous graph expansion, including the exact unchanged host attachment coordinate.
For the proof, the preceding component lemmas are composed or specialized at the stated parameters.
\lean{Copying/CopyingTemplateInputAssembly}{12}{23}{CopyingTemplate.expanded_entry}
\par\penalty0

No original filling outside the entered component is discarded. The constructed residual has the actual host token and exact two-transition prefixes on both sides, and sharing is used only for its proved correction.
For the proof, the defined maps and incidence equations are substituted.
\lean{Copying/CopyingTemplateInputAssembly}{28}{84}{CopyingTemplate.entry_diamond}
\par\penalty0

The remaining local obligation for multi-input copying: an actual request/compute prefix on each side and a certified original residual.
\lean{Copying/CopyingTemplateInputAssembly}{88}{95}{CopyingTemplate.OtherEntries}
\par\penalty0

Case II handles the distinguished input with its exact 2/0 costs. Case III handles every other input with exact 2/2 costs and the independently completed sharing correction. The local diagrams, rather than a semantic equivalence assumption, are the only argument to this assembly theorem.
For the proof, constructor and boundary cases are separated explicitly; the remaining natural-number index bounds follow by linear arithmetic; the defined maps and incidence equations are substituted.
\lean{Copying/CopyingTemplateInputAssembly}{101}{127}{CopyingTemplate.inputSafe_of_otherEntries}
\par\penalty0

\subsection{Copying Template Input Safety}
Unconditional input safety of the original unrestricted copying template. The distinguished input consumes the literal copy in two transitions. Every other ordered input uses two actual copy prefixes and the proved two-copy sharing interchange, retaining all original payloads and residual trees.

Each non-distinguished input of the literal family has its actual 2/2 entry diamond and an original copying residual corrected by sharing alone.
For the proof, constructor and boundary cases are separated explicitly; the remaining natural-number index bounds follow by linear arithmetic; the defined maps and incidence equations are substituted.
\lean{Copying/CopyingTemplateInputSafety}{15}{35}{CopyingTemplate.Literal.other_entry}
\par\penalty0

Independent representations of both endpoints preserve the literal local entry diagrams. No determinism, residual, or semantic premise remains.
For the proof, constructor and boundary cases are separated explicitly.
\lean{Copying/CopyingTemplateInputSafety}{39}{48}{CopyingTemplate.otherEntries}
\par\penalty0

Original copying is input safe in every focus-free surrounding context for every cost comparison permitting its exact distinguished 2/0 prefix.
For the proof, the preceding component lemmas are composed or specialized at the stated parameters.
\lean{Copying/CopyingTemplateInputSafety}{52}{54}{CopyingTemplate.template_inputSafe_cost}
\par\penalty0

Forgetting costs gives unrestricted copying input safety directly.
For the proof, the preceding component lemmas are composed or specialized at the stated parameters.
\lean{Copying/CopyingTemplateInputSafety}{57}{59}{CopyingTemplate.template_inputSafe}
\par\penalty0

The paper's quantitative input-safety comparison retains the original left-to-right improvement rather than erasing the distinguished copying cost.
For the proof, the preceding component lemmas are composed or specialized at the stated parameters.
\lean{Copying/CopyingTemplateInputSafety}{63}{65}{CopyingTemplate.template_inputSafe_improvement}
\par\penalty0

\subsection{Copying Template Laws}
Output closure and representation padding for the literal copying pre-template, retaining its complete multi-input shape.

A distributor contains only sharing labels. Hence no operation-only path can arrive at any of its output wires.
For the proof, constructor and boundary cases are separated explicitly; the defined maps and incidence equations are substituted.
\lean{Copying/CopyingTemplateLaws}{12}{27}{CopyingTemplate.distributor_oneWay}
\par\penalty0

The literal right side is one-way because its final term and atom distributors intercept all its outputs, regardless of the payload body.
For the proof, constructor and boundary cases are separated explicitly.
\lean{Copying/CopyingTemplateLaws}{31}{35}{CopyingTemplate.Literal.outputClosed}
\par\penalty0

The output Closed assertion computes the indicated ordered boundary coordinate in the assembled graph.
For the proof, constructor and boundary cases are separated explicitly.
\lean{Copying/CopyingTemplateLaws}{37}{40}{CopyingTemplate.Pair.outputClosed}
\par\penalty0

The template output Closed assertion computes the indicated ordered boundary coordinate in the assembled graph.
For the proof, the preceding component lemmas are composed or specialized at the stated parameters.
\lean{Copying/CopyingTemplateLaws}{42}{42}{CopyingTemplate.template_outputClosed}
\par\penalty0

Raising the allowance preserves the same tree and raises the original payload recursively; both sides remain the original copying construction.
For the proof, constructor and boundary cases are separated explicitly.
\lean{Copying/CopyingTemplateLaws}{46}{52}{CopyingTemplate.Literal.raise}
\par\penalty0

The raise assertion is the compatibility law for the indicated component of this construction. It relates the actual finite coordinates or certificates, with the hypotheses stated in the linked declaration.
For the proof, constructor and boundary cases are separated explicitly.
\lean{Copying/CopyingTemplateLaws}{54}{60}{CopyingTemplate.Pair.raise}
\par\penalty0

The template padding Closed assertion is the compatibility law for the indicated component of this construction. It relates the actual finite coordinates or certificates, with the hypotheses stated in the linked declaration.
For the proof, the preceding component lemmas are composed or specialized at the stated parameters.
\lean{Copying/CopyingTemplateLaws}{62}{62}{CopyingTemplate.template_paddingClosed}
\par\penalty0

\subsection{Copying Template Residual}
A literal copying residual followed by an independently proved exact sharing-tree correction. The original copying relation is not enlarged by transitivity or by an assumed contextual equivalence.

Copying retains every payload label on both sides, and adds only sharing labels; therefore the literal pair preserves recursive hole-freedom exactly.
For the proof, constructor and boundary cases are separated explicitly; the defined maps and incidence equations are substituted.
\lean{Copying/CopyingTemplateResidual}{14}{34}{CopyingTemplate.Literal.holeFree}
\par\penalty0

The hole Free assertion excludes unfilled holes throughout the indicated graph and every nested box.
For the proof, constructor and boundary cases are separated explicitly.
\lean{Copying/CopyingTemplateResidual}{36}{40}{CopyingTemplate.Pair.holeFree}
\par\penalty0

The first link is an actual recursive copying filling. Only the second link uses the completed contraction-tree algebra.
\lean{Copying/CopyingTemplateResidual}{44}{49}{CopyingTemplate.Residual}
\par\penalty0

This defines the direct data for the construction just described. Boundary positions are kept distinct from newly inserted internal vertices.
\lean{Copying/CopyingTemplateResidual}{52}{54}{CopyingTemplate.Residual.direct}
\par\penalty0

The congr map is an ordered recursive graph isomorphism between the stated original constructions.
\lean{Copying/CopyingTemplateResidual}{56}{63}{CopyingTemplate.Residual.congr}
\par\penalty0

The same original copying specimen is a quasi-specimen up to exact execution equivalence. In particular the common residual's right root is derived from the actual right endpoint through proved sharing reflection.
For the proof, the defined maps and incidence equations are substituted.
\lean{Copying/CopyingTemplateResidual}{68}{86}{CopyingTemplate.Residual.quasi_equiv}
\par\penalty0

\subsection{Copying Template Residual Expansion}
Closure of the two-link copying residual under arbitrary simultaneous expansions. Original copying contributions remain in the first link.

Untouched components keep their original certificates in the copying link, and remain literal reflexive components during the sharing correction.
\lean{Copying/CopyingTemplateResidualExpansion}{14}{34}{CopyingTemplate.Residual.expandSelected}
\par\penalty0

A selected actual rewrite and its two-link residual lift together through a whole simultaneous expansion. The complement is the original derivation's finite frame, with every other contribution explicitly retained.
For the proof, the preceding component lemmas are composed or specialized at the stated parameters.
\lean{Copying/CopyingTemplateResidualExpansion}{39}{61}{CopyingTemplate.Residual.assemble}
\par\penalty0

\subsection{Parametricity Entry}
Every input of Figure 20 reaches exactly a surface lambda head and returns immediately. No execution is allowed to enter a private box merely by probing its value at the original template boundary.

This defines the Lambda Inputs data for the construction just described. Boundary positions are kept distinct from newly inserted internal vertices.
\lean{Copying/ParametricityEntry}{11}{12}{Parametricity.LambdaInputs}
\par\penalty0

The lambda Inputs payload assertion computes the indicated ordered boundary coordinate in the assembled graph.
For the proof, the remaining natural-number index bounds follow by linear arithmetic; the defined maps and incidence equations are substituted.
\lean{Copying/ParametricityEntry}{14}{21}{Parametricity.lambdaInputs_payload}
\par\penalty0

Retaining a serial suffix preserves every source lambda head, including when its output is attached to the private sharing tree.
For the proof, constructor and boundary cases are separated explicitly; the defined maps and incidence equations are substituted.
\lean{Copying/ParametricityEntry}{25}{34}{Parametricity.LambdaInputs.serial}
\par\penalty0

Ordered disjoint juxtaposition preserves the surface lambda-head property at every actual input, without inserting forwarding edges.
For the proof, constructor and boundary cases are separated explicitly; the remaining natural-number index bounds follow by linear arithmetic; the defined maps and incidence equations are substituted.
\lean{Copying/ParametricityEntry}{38}{55}{Parametricity.LambdaInputs.tensor}
\par\penalty0

Equality transport changes no input coordinate or edge occurrence.
For the proof, constructor and boundary cases are separated explicitly.
\lean{Copying/ParametricityEntry}{58}{61}{Parametricity.LambdaInputs.reboundary}
\par\penalty0

The copies lambda Inputs assertion computes the indicated ordered boundary coordinate in the assembled graph.
For the proof, induction follows the original derivation or finite structure.
\lean{Copying/ParametricityEntry}{63}{67}{Parametricity.copies_lambdaInputs}
\par\penalty0

The left lambda Inputs assertion computes the indicated ordered boundary coordinate in the assembled graph.
For the proof, the preceding component lemmas are composed or specialized at the stated parameters.
\lean{Copying/ParametricityEntry}{69}{72}{Parametricity.left_lambdaInputs}
\par\penalty0

The right lambda Inputs assertion computes the indicated ordered boundary coordinate in the assembled graph.
For the proof, the preceding component lemmas are composed or specialized at the stated parameters.
\lean{Copying/ParametricityEntry}{74}{76}{Parametricity.right_lambdaInputs}
\par\penalty0

The original search transition immediately returns from every surface lambda value; no callback or private-store action is executed at entry.
For the proof, the defined maps and incidence equations are substituted.
\lean{Copying/ParametricityEntry}{80}{83}{Parametricity.LambdaInputs.returns}
\par\penalty0

\subsection{Parametricity Graphs}
The literal schematic graphs of Figure 20. The two abstraction families are retained separately and in their displayed order. The private reference is shared through an arbitrary TERM contraction tree, followed by the single instance, canonical ATOM D1,1, store, and numeral 1. These graphs are not a replacement definition of the inductive term translation.

The ignored w is the first bound output; dereferencing reaches the sole free variable x, exactly as in the smaller left abstraction of Figure 20.
\lean{Copying/ParametricityGraphs}{17}{19}{Parametricity.readBody}
\par\penalty0

The smaller right abstraction ignores w and returns literal numeral 1.
\lean{Copying/ParametricityGraphs}{22}{24}{Parametricity.oneBody}
\par\penalty0

This defines the small Left data for the construction just described. Boundary positions are kept distinct from newly inserted internal vertices.
\lean{Copying/ParametricityGraphs}{26}{26}{Parametricity.smallLeft}
\par\penalty0

This defines the small Right data for the construction just described. Boundary positions are kept distinct from newly inserted internal vertices.
\lean{Copying/ParametricityGraphs}{27}{27}{Parametricity.smallRight}
\par\penalty0

The callback f() retains the original displayed single-use term contraction tree D1,1 and the literal unit edge.
\lean{Copying/ParametricityGraphs}{31}{34}{Parametricity.callback}
\par\penalty0

Swap only the two equal-typed output positions, making f the bound output and x the free output of the larger left abstraction.
\lean{Copying/ParametricityGraphs}{38}{42}{Parametricity.boundFirst}
\par\penalty0

This defines the call Read Body data for the construction just described. Boundary positions are kept distinct from newly inserted internal vertices.
\lean{Copying/ParametricityGraphs}{44}{47}{Parametricity.callReadBody}
\par\penalty0

This defines the call One Body data for the construction just described. Boundary positions are kept distinct from newly inserted internal vertices.
\lean{Copying/ParametricityGraphs}{49}{51}{Parametricity.callOneBody}
\par\penalty0

This defines the large Left data for the construction just described. Boundary positions are kept distinct from newly inserted internal vertices.
\lean{Copying/ParametricityGraphs}{53}{53}{Parametricity.largeLeft}
\par\penalty0

This defines the large Right data for the construction just described. Boundary positions are kept distinct from newly inserted internal vertices.
\lean{Copying/ParametricityGraphs}{54}{54}{Parametricity.largeRight}
\par\penalty0

Ordered disjoint copies, including zero copies. No identity edge is inserted: the zero case is the actual empty hypernet.
\lean{Copying/ParametricityGraphs}{58}{63}{Parametricity.copies}
\par\penalty0

m1 larger abstractions precede m2 smaller abstractions; all free x outputs are kept, in that same order, for the private sharing tree.
\lean{Copying/ParametricityGraphs}{67}{72}{Parametricity.leftFamily}
\par\penalty0

This defines the right Family data for the construction just described. Boundary positions are kept distinct from newly inserted internal vertices.
\lean{Copying/ParametricityGraphs}{74}{78}{Parametricity.rightFamily}
\par\penalty0

The private reference contains the canonical atom D1,1; it is not replaced by a direct wire or a simplified store encoding.
\lean{Copying/ParametricityGraphs}{82}{83}{Parametricity.privateCell}
\par\penalty0

The entire finite family alone owns every input of C and the unique instance feeding its private store. There is no free resource output.
\lean{Copying/ParametricityGraphs}{87}{89}{Parametricity.left}
\par\penalty0

This defines the right graph fragment for the construction just described. Boundary positions are kept distinct from newly inserted internal vertices.
\lean{Copying/ParametricityGraphs}{91}{92}{Parametricity.right}
\par\penalty0

\subsection{Parametricity Labels}
Recursive label invariants of the literal Figure 20 graphs, retaining every weakening, contraction, box body and private store edge.

Exactly the primitive constructors used by the two schematic families and their shared private cell. This is a syntactic invariant, not a semantic assumption about arbitrary callbacks supplied by observing contexts.
\lean{Copying/ParametricityLabels}{12}{16}{Parametricity.PrimitiveLabel}
\par\penalty0

The read Body labels assertion verifies the indicated label compatibility or label exclusion at the original coordinate.
For the proof, the preceding component lemmas are composed or specialized at the stated parameters.
\lean{Copying/ParametricityLabels}{26}{27}{Parametricity.readBody_labels}
\par\penalty0

The one Body labels assertion verifies the indicated label compatibility or label exclusion at the original coordinate.
For the proof, the preceding component lemmas are composed or specialized at the stated parameters.
\lean{Copying/ParametricityLabels}{29}{30}{Parametricity.oneBody_labels}
\par\penalty0

The small Left labels assertion verifies the indicated label compatibility or label exclusion at the original coordinate.
For the proof, the preceding component lemmas are composed or specialized at the stated parameters.
\lean{Copying/ParametricityLabels}{32}{32}{Parametricity.smallLeft_labels}
\par\penalty0

The small Right labels assertion verifies the indicated label compatibility or label exclusion at the original coordinate.
For the proof, the preceding component lemmas are composed or specialized at the stated parameters.
\lean{Copying/ParametricityLabels}{33}{33}{Parametricity.smallRight_labels}
\par\penalty0

The callback labels assertion verifies the indicated label compatibility or label exclusion at the original coordinate.
For the proof, the defined maps and incidence equations are substituted.
\lean{Copying/ParametricityLabels}{35}{44}{Parametricity.callback_labels}
\par\penalty0

The call Read Body labels assertion verifies the indicated label compatibility or label exclusion at the original coordinate.
For the proof, the preceding component lemmas are composed or specialized at the stated parameters.
\lean{Copying/ParametricityLabels}{46}{48}{Parametricity.callReadBody_labels}
\par\penalty0

The call One Body labels assertion verifies the indicated label compatibility or label exclusion at the original coordinate.
For the proof, the preceding component lemmas are composed or specialized at the stated parameters.
\lean{Copying/ParametricityLabels}{50}{52}{Parametricity.callOneBody_labels}
\par\penalty0

The large Left labels assertion verifies the indicated label compatibility or label exclusion at the original coordinate.
For the proof, the preceding component lemmas are composed or specialized at the stated parameters.
\lean{Copying/ParametricityLabels}{54}{54}{Parametricity.largeLeft_labels}
\par\penalty0

The large Right labels assertion verifies the indicated label compatibility or label exclusion at the original coordinate.
For the proof, the preceding component lemmas are composed or specialized at the stated parameters.
\lean{Copying/ParametricityLabels}{55}{55}{Parametricity.largeRight_labels}
\par\penalty0

Replication is finite disjoint assembly; it introduces no primitive.
For the proof, induction follows the original derivation or finite structure.
\lean{Copying/ParametricityLabels}{58}{62}{Parametricity.copies_labels}
\par\penalty0

The left Family labels assertion verifies the indicated label compatibility or label exclusion at the original coordinate.
For the proof, the preceding component lemmas are composed or specialized at the stated parameters.
\lean{Copying/ParametricityLabels}{64}{66}{Parametricity.leftFamily_labels}
\par\penalty0

The right Family labels assertion verifies the indicated label compatibility or label exclusion at the original coordinate.
For the proof, the preceding component lemmas are composed or specialized at the stated parameters.
\lean{Copying/ParametricityLabels}{68}{70}{Parametricity.rightFamily_labels}
\par\penalty0

The private Cell labels assertion verifies the indicated label compatibility or label exclusion at the original coordinate.
For the proof, the defined maps and incidence equations are substituted.
\lean{Copying/ParametricityLabels}{72}{83}{Parametricity.privateCell_labels}
\par\penalty0

The left labels assertion verifies the indicated label compatibility or label exclusion at the original coordinate.
For the proof, constructor and boundary cases are separated explicitly; the defined maps and incidence equations are substituted.
\lean{Copying/ParametricityLabels}{85}{93}{Parametricity.left_labels}
\par\penalty0

The right labels assertion verifies the indicated label compatibility or label exclusion at the original coordinate.
For the proof, the preceding component lemmas are composed or specialized at the stated parameters.
\lean{Copying/ParametricityLabels}{95}{96}{Parametricity.right_labels}
\par\penalty0

The full recursively nested figure is both focus-free and hole-free.
For the proof, induction follows the original derivation or finite structure; constructor and boundary cases are separated explicitly.
\lean{Copying/ParametricityLabels}{99}{118}{Parametricity.labels_closed}
\par\penalty0

\subsection{Parametricity Opening}
The actual original beta steps for Figure 20's smaller abstractions. Both branches retain the entire unused stable argument under its weakening; only the live dereference versus literal numeral distinguishes the reducts.

Opening preserves the literal bound weakening and the private read.
For the proof, the defined maps and incidence equations are substituted.
\lean{Copying/ParametricityOpening}{13}{16}{Parametricity.open_readBody}
\par\penalty0

Opening preserves the literal bound weakening and numeral 1.
For the proof, the defined maps and incidence equations are substituted.
\lean{Copying/ParametricityOpening}{19}{22}{Parametricity.open_oneBody}
\par\penalty0

The exact first reduct on the private-read branch, including all original resource outputs of the unused argument.
\lean{Copying/ParametricityOpening}{26}{28}{Parametricity.readOpened}
\par\penalty0

The exact first reduct on the constant branch. The same discarded argument remains as a separate component and cannot be erased by beta.
\lean{Copying/ParametricityOpening}{32}{34}{Parametricity.oneOpened}
\par\penalty0

The factorization is an actual recursive graph isomorphism, not a semantic argument-discarding rule.
\lean{Copying/ParametricityOpening}{38}{42}{Parametricity.read_openedIso}
\par\penalty0

The one opened Iso map is an ordered recursive graph isomorphism between the stated original constructions.
\lean{Copying/ParametricityOpening}{44}{48}{Parametricity.one_openedIso}
\par\penalty0

The left branch takes precisely the original stable beta rule, with the same active input before and after the step.
For the proof, the preceding component lemmas are composed or specialized at the stated parameters.
\lean{Copying/ParametricityOpening}{52}{56}{Parametricity.read_beta}
\par\penalty0

The right branch likewise performs its actual beta computation.
For the proof, the preceding component lemmas are composed or specialized at the stated parameters.
\lean{Copying/ParametricityOpening}{59}{63}{Parametricity.one_beta}
\par\penalty0

\subsection{Parametricity Padding}
Literal Figure 20 is closed under the original depth padding. Every ordered family, private sharing tree, distributor, box and store is retained.

The read Body raise assertion is the compatibility law for the indicated component of this construction. It relates the actual finite coordinates or certificates, with the hypotheses stated in the linked declaration.
For the proof, the defined maps and incidence equations are substituted.
\lean{Copying/ParametricityPadding}{15}{17}{Parametricity.readBody_raise}
\par\penalty0

The one Body raise assertion is the compatibility law for the indicated component of this construction. It relates the actual finite coordinates or certificates, with the hypotheses stated in the linked declaration.
For the proof, the defined maps and incidence equations are substituted.
\lean{Copying/ParametricityPadding}{19}{21}{Parametricity.oneBody_raise}
\par\penalty0

The small Left raise map is an ordered recursive graph isomorphism between the stated original constructions.
\lean{Copying/ParametricityPadding}{23}{24}{Parametricity.smallLeft_raise}
\par\penalty0

The small Right raise map is an ordered recursive graph isomorphism between the stated original constructions.
\lean{Copying/ParametricityPadding}{26}{27}{Parametricity.smallRight_raise}
\par\penalty0

The callback raise map is an ordered recursive graph isomorphism between the stated original constructions.
\lean{Copying/ParametricityPadding}{29}{36}{Parametricity.callback_raise}
\par\penalty0

The call Read Body raise map is an ordered recursive graph isomorphism between the stated original constructions.
\lean{Copying/ParametricityPadding}{38}{47}{Parametricity.callReadBody_raise}
\par\penalty0

The call One Body raise map is an ordered recursive graph isomorphism between the stated original constructions.
\lean{Copying/ParametricityPadding}{49}{55}{Parametricity.callOneBody_raise}
\par\penalty0

The large Left raise map is an ordered recursive graph isomorphism between the stated original constructions.
\lean{Copying/ParametricityPadding}{57}{58}{Parametricity.largeLeft_raise}
\par\penalty0

The large Right raise map is an ordered recursive graph isomorphism between the stated original constructions.
\lean{Copying/ParametricityPadding}{60}{61}{Parametricity.largeRight_raise}
\par\penalty0

The copies congr map is an ordered recursive graph isomorphism between the stated original constructions.
\lean{Copying/ParametricityPadding}{63}{67}{Parametricity.copies_congr}
\par\penalty0

The copies raise map is an ordered recursive graph isomorphism between the stated original constructions.
\lean{Copying/ParametricityPadding}{69}{84}{Parametricity.copies_raise}
\par\penalty0

The left Family raise map is an ordered recursive graph isomorphism between the stated original constructions.
\lean{Copying/ParametricityPadding}{86}{92}{Parametricity.leftFamily_raise}
\par\penalty0

The right Family raise map is an ordered recursive graph isomorphism between the stated original constructions.
\lean{Copying/ParametricityPadding}{94}{100}{Parametricity.rightFamily_raise}
\par\penalty0

The private Cell raise map is an ordered recursive graph isomorphism between the stated original constructions.
\lean{Copying/ParametricityPadding}{102}{118}{Parametricity.privateCell_raise}
\par\penalty0

The left raise map is an ordered recursive graph isomorphism between the stated original constructions.
\lean{Copying/ParametricityPadding}{120}{126}{Parametricity.left_raise}
\par\penalty0

The raise assertion is the compatibility law for the indicated component of this construction. It relates the actual finite coordinates or certificates, with the hypotheses stated in the linked declaration.
For the proof, constructor and boundary cases are separated explicitly.
\lean{Copying/ParametricityPadding}{128}{135}{Parametricity.Pair.raise}
\par\penalty0

The template padding Closed assertion is the compatibility law for the indicated component of this construction. It relates the actual finite coordinates or certificates, with the hypotheses stated in the linked declaration.
For the proof, the preceding component lemmas are composed or specialized at the stated parameters.
\lean{Copying/ParametricityPadding}{137}{137}{Parametricity.template_paddingClosed}
\par\penalty0

\subsection{Parametricity Private Application}
The smaller private-reading abstraction's original beta step and exact nine-transition prefix, with an arbitrary stable argument and its literal unused resources retained in parallel.

This defines the cell data for the construction just described. Boundary positions are kept distinct from newly inserted internal vertices.
\lean{Copying/ParametricityPrivateApplication}{13}{13}{Parametricity.PrivateApplication.cell}
\par\penalty0

This defines the discarded data for the construction just described. Boundary positions are kept distinct from newly inserted internal vertices.
\lean{Copying/ParametricityPrivateApplication}{14}{14}{Parametricity.PrivateApplication.discarded}
\par\penalty0

This defines the opened data for the construction just described. Boundary positions are kept distinct from newly inserted internal vertices.
\lean{Copying/ParametricityPrivateApplication}{16}{17}{Parametricity.PrivateApplication.opened}
\par\penalty0

The private Occurrence map is a complete occurrence, preserving original labels and every ordered source and target attachment.
\lean{Copying/ParametricityPrivateApplication}{19}{21}{Parametricity.PrivateApplication.privateOccurrence}
\par\penalty0

The argument Occurrence map is a complete occurrence, preserving original labels and every ordered source and target attachment.
\lean{Copying/ParametricityPrivateApplication}{23}{25}{Parametricity.PrivateApplication.argumentOccurrence}
\par\penalty0

This defines the reader data for the construction just described. Boundary positions are kept distinct from newly inserted internal vertices.
\lean{Copying/ParametricityPrivateApplication}{27}{27}{Parametricity.PrivateApplication.reader}
\par\penalty0

The reader label assertion verifies the indicated label compatibility or label exclusion at the original coordinate.
Both sides are definitionally equal.
\lean{Copying/ParametricityPrivateApplication}{28}{29}{Parametricity.PrivateApplication.reader_label}
\par\penalty0

The reader Occurrence map is a complete occurrence, preserving original labels and every ordered source and target attachment.
\lean{Copying/ParametricityPrivateApplication}{31}{32}{Parametricity.PrivateApplication.readerOccurrence}
\par\penalty0

This defines the active data for the construction just described. Boundary positions are kept distinct from newly inserted internal vertices.
\lean{Copying/ParametricityPrivateApplication}{34}{35}{Parametricity.PrivateApplication.active}
\par\penalty0

The reader input assertion computes the indicated ordered boundary coordinate in the assembled graph.
For the proof, the defined maps and incidence equations are substituted.
\lean{Copying/ParametricityPrivateApplication}{37}{46}{Parametricity.PrivateApplication.reader_input}
\par\penalty0

The components map is a complete occurrence, preserving original labels and every ordered source and target attachment.
\lean{Copying/ParametricityPrivateApplication}{48}{52}{Parametricity.PrivateApplication.components}
\par\penalty0

The disjoint assertion excludes a shared edge or coordinate between the indicated independently inserted components.
For the proof, constructor and boundary cases are separated explicitly; the defined maps and incidence equations are substituted.
\lean{Copying/ParametricityPrivateApplication}{54}{61}{Parametricity.PrivateApplication.disjoint}
\par\penalty0

The opened original read-body redex has its complete original occurrence in the actual private tree/store graph plus the unchanged discarded argument.
\lean{Copying/ParametricityPrivateApplication}{65}{69}{Parametricity.PrivateApplication.bodyOccurrence}
\par\penalty0

The body input assertion computes the indicated ordered boundary coordinate in the assembled graph.
For the proof, the defined maps and incidence equations are substituted.
\lean{Copying/ParametricityPrivateApplication}{71}{73}{Parametricity.PrivateApplication.body_input}
\par\penalty0

This source is the actual original beta redex in its computed private-cell context, preserving the stable argument and all its resource outputs.
\lean{Copying/ParametricityPrivateApplication}{77}{79}{Parametricity.PrivateApplication.before}
\par\penalty0

This defines the token data for the construction just described. Boundary positions are kept distinct from newly inserted internal vertices.
\lean{Copying/ParametricityPrivateApplication}{81}{83}{Parametricity.PrivateApplication.token}
\par\penalty0

The beta assertion is the compatibility law for the indicated component of this construction. It relates the actual finite coordinates or certificates, with the hypotheses stated in the linked declaration.
For the proof, the defined maps and incidence equations are substituted.
\lean{Copying/ParametricityPrivateApplication}{85}{90}{Parametricity.PrivateApplication.beta}
\par\penalty0

Nine actual transitions: original beta, then the exact eight-transition private readback. The endpoint records its unchanged complementary resources.
For the proof, the preceding component lemmas are composed or specialized at the stated parameters.
\lean{Copying/ParametricityPrivateApplication}{94}{102}{Parametricity.PrivateApplication.run}
\par\penalty0

\subsection{Parametricity Private Execution}
The literal private-cell readback performs eight transitions after its reader becomes active: five expose the store read, and three read and copy 1. All unmatched references and weakening branches remain in the computed cut.

The one label assertion verifies the indicated label compatibility or label exclusion at the original coordinate.
For the proof, constructor and boundary cases are separated explicitly; the defined maps and incidence equations are substituted.
\lean{Copying/ParametricityPrivateExecution}{12}{13}{Parametricity.PrivateRead.one_label}
\par\penalty0

The numeral Occurrence map is a complete occurrence, preserving original labels and every ordered source and target attachment.
\lean{Copying/ParametricityPrivateExecution}{15}{16}{Parametricity.PrivateRead.numeralOccurrence}
\par\penalty0

The read token assertion is the compatibility law for the indicated component of this construction. It relates the actual finite coordinates or certificates, with the hypotheses stated in the linked declaration.
Both sides are definitionally equal.
\lean{Copying/ParametricityPrivateExecution}{18}{19}{Parametricity.PrivateRead.read_token}
\par\penalty0

The content boundary assertion is the compatibility law for the indicated component of this construction. It relates the actual finite coordinates or certificates, with the hypotheses stated in the linked declaration.
For the proof, constructor and boundary cases are separated explicitly; the defined maps and incidence equations are substituted.
\lean{Copying/ParametricityPrivateExecution}{21}{28}{Parametricity.PrivateRead.content_boundary}
\par\penalty0

The read numeral disjoint assertion excludes a shared edge or coordinate between the indicated independently inserted components.
For the proof, constructor and boundary cases are separated explicitly; the defined maps and incidence equations are substituted.
\lean{Copying/ParametricityPrivateExecution}{30}{41}{Parametricity.PrivateRead.read_numeral_disjoint}
\par\penalty0

The selected local read and the actual stored numeral form one complete serial occurrence; the old private atom and every inactive branch are retained.
\lean{Copying/ParametricityPrivateExecution}{45}{48}{Parametricity.PrivateRead.tailOccurrence}
\par\penalty0

The tail token assertion is the compatibility law for the indicated component of this construction. It relates the actual finite coordinates or certificates, with the hypotheses stated in the linked declaration.
For the proof, the defined maps and incidence equations are substituted.
\lean{Copying/ParametricityPrivateExecution}{50}{54}{Parametricity.PrivateRead.tail_token}
\par\penalty0

This defines the after data for the construction just described. Boundary positions are kept distinct from newly inserted internal vertices.
\lean{Copying/ParametricityPrivateExecution}{56}{59}{Parametricity.PrivateRead.after}
\par\penalty0

The real reader-ready graph performs the paper's final three transitions at its actual token, ending in the explicit untouched-complement replacement.
For the proof, constructor and boundary cases are separated explicitly; the defined maps and incidence equations are substituted.
\lean{Copying/ParametricityPrivateExecution}{63}{94}{Parametricity.PrivateRead.readback}
\par\penalty0

Eight literal transitions after opening the smaller private-reading closure, stopping before the shared passive return of the copied numeral.
For the proof, the defined maps and incidence equations are substituted.
\lean{Copying/ParametricityPrivateExecution}{98}{103}{Parametricity.PrivateRead.run}
\par\penalty0

\subsection{Parametricity Private Read}
The private store's complete read match after duplicating its instance. All three actual atom contractions are selected, while their unused sibling attachments remain in the original complement.

This defines the tail data for the construction just described. Boundary positions are kept distinct from newly inserted internal vertices.
\lean{Copying/ParametricityPrivateRead}{13}{14}{Parametricity.PrivateRead.tail}
\par\penalty0

This defines the before data for the construction just described. Boundary positions are kept distinct from newly inserted internal vertices.
\lean{Copying/ParametricityPrivateRead}{16}{17}{Parametricity.PrivateRead.before}
\par\penalty0

This defines the ready data for the construction just described. Boundary positions are kept distinct from newly inserted internal vertices.
\lean{Copying/ParametricityPrivateRead}{18}{19}{Parametricity.PrivateRead.ready}
\par\penalty0

This defines the reader data for the construction just described. Boundary positions are kept distinct from newly inserted internal vertices.
\lean{Copying/ParametricityPrivateRead}{21}{21}{Parametricity.PrivateRead.reader}
\par\penalty0

This defines the copy Edge data for the construction just described. Boundary positions are kept distinct from newly inserted internal vertices.
\lean{Copying/ParametricityPrivateRead}{22}{23}{Parametricity.PrivateRead.copyEdge}
\par\penalty0

This defines the instance Edge data for the construction just described. Boundary positions are kept distinct from newly inserted internal vertices.
\lean{Copying/ParametricityPrivateRead}{24}{25}{Parametricity.PrivateRead.instanceEdge}
\par\penalty0

This defines the new Share data for the construction just described. Boundary positions are kept distinct from newly inserted internal vertices.
\lean{Copying/ParametricityPrivateRead}{26}{27}{Parametricity.PrivateRead.newShare}
\par\penalty0

This defines the old Share data for the construction just described. Boundary positions are kept distinct from newly inserted internal vertices.
\lean{Copying/ParametricityPrivateRead}{28}{29}{Parametricity.PrivateRead.oldShare}
\par\penalty0

This defines the store data for the construction just described. Boundary positions are kept distinct from newly inserted internal vertices.
\lean{Copying/ParametricityPrivateRead}{30}{31}{Parametricity.PrivateRead.store}
\par\penalty0

This defines the one data for the construction just described. Boundary positions are kept distinct from newly inserted internal vertices.
\lean{Copying/ParametricityPrivateRead}{32}{33}{Parametricity.PrivateRead.one}
\par\penalty0

This defines the atom Edges data for the construction just described. Boundary positions are kept distinct from newly inserted internal vertices.
\lean{Copying/ParametricityPrivateRead}{35}{36}{Parametricity.PrivateRead.atomEdges}
\par\penalty0

The atom label assertion verifies the indicated label compatibility or label exclusion at the original coordinate.
For the proof, constructor and boundary cases are separated explicitly; the remaining natural-number index bounds follow by linear arithmetic; the defined maps and incidence equations are substituted.
\lean{Copying/ParametricityPrivateRead}{38}{41}{Parametricity.PrivateRead.atom_label}
\par\penalty0

This defines the atom Path data for the construction just described. Boundary positions are kept distinct from newly inserted internal vertices.
\lean{Copying/ParametricityPrivateRead}{43}{47}{Parametricity.PrivateRead.atomPath}
\par\penalty0

This defines the atom Sides data for the construction just described. Boundary positions are kept distinct from newly inserted internal vertices.
\lean{Copying/ParametricityPrivateRead}{48}{50}{Parametricity.PrivateRead.atomSides}
\par\penalty0

This defines the orientation data for the construction just described. Boundary positions are kept distinct from newly inserted internal vertices.
\lean{Copying/ParametricityPrivateRead}{51}{51}{Parametricity.PrivateRead.orientation}
\par\penalty0

The observed source path crosses the new atom distributor and then the original D1,1 at the allocated cell; no atom wire is contracted away.
For the proof, constructor and boundary cases are separated explicitly; the remaining natural-number index bounds follow by linear arithmetic; the defined maps and incidence equations are substituted.
\lean{Copying/ParametricityPrivateRead}{55}{60}{Parametricity.PrivateRead.atom_source}
\par\penalty0

The atom target identity gives the entire ordered target list after applying the construction, with output order retained.
For the proof, constructor and boundary cases are separated explicitly; the remaining natural-number index bounds follow by linear arithmetic; the defined maps and incidence equations are substituted.
\lean{Copying/ParametricityPrivateRead}{62}{65}{Parametricity.PrivateRead.atom_target}
\par\penalty0

The atom injective assertion states injectivity of the indicated coordinate map: equal image coordinates have equal original coordinates.
For the proof, constructor and boundary cases are separated explicitly; the remaining natural-number index bounds follow by linear arithmetic; the defined maps and incidence equations are substituted.
\lean{Copying/ParametricityPrivateRead}{67}{81}{Parametricity.PrivateRead.atom_injective}
\par\penalty0

The tree Occurrence map is a complete occurrence, preserving original labels and every ordered source and target attachment.
\lean{Copying/ParametricityPrivateRead}{83}{87}{Parametricity.PrivateRead.treeOccurrence}
\par\penalty0

This defines the argument data for the construction just described. Boundary positions are kept distinct from newly inserted internal vertices.
\lean{Copying/ParametricityPrivateRead}{89}{90}{Parametricity.PrivateRead.argument}
\par\penalty0

This defines the store Value data for the construction just described. Boundary positions are kept distinct from newly inserted internal vertices.
\lean{Copying/ParametricityPrivateRead}{91}{92}{Parametricity.PrivateRead.storeValue}
\par\penalty0

The reader label assertion verifies the indicated label compatibility or label exclusion at the original coordinate.
For the proof, constructor and boundary cases are separated explicitly; the defined maps and incidence equations are substituted.
\lean{Copying/ParametricityPrivateRead}{94}{95}{Parametricity.PrivateRead.reader_label}
\par\penalty0

The instance label assertion verifies the indicated label compatibility or label exclusion at the original coordinate.
For the proof, constructor and boundary cases are separated explicitly; the defined maps and incidence equations are substituted.
\lean{Copying/ParametricityPrivateRead}{96}{97}{Parametricity.PrivateRead.instance_label}
\par\penalty0

The store label assertion verifies the indicated label compatibility or label exclusion at the original coordinate.
For the proof, constructor and boundary cases are separated explicitly; the defined maps and incidence equations are substituted.
\lean{Copying/ParametricityPrivateRead}{98}{98}{Parametricity.PrivateRead.store_label}
\par\penalty0

The reader source identity gives the entire ordered source list after applying the construction, rather than only its underlying set.
For the proof, constructor and boundary cases are separated explicitly; the defined maps and incidence equations are substituted.
\lean{Copying/ParametricityPrivateRead}{100}{113}{Parametricity.PrivateRead.reader_source}
\par\penalty0

The reader target identity gives the entire ordered target list after applying the construction, with output order retained.
For the proof, constructor and boundary cases are separated explicitly; the defined maps and incidence equations are substituted.
\lean{Copying/ParametricityPrivateRead}{114}{114}{Parametricity.PrivateRead.reader_target}
\par\penalty0

The instance source identity gives the entire ordered source list after applying the construction, rather than only its underlying set.
For the proof, constructor and boundary cases are separated explicitly; the defined maps and incidence equations are substituted.
\lean{Copying/ParametricityPrivateRead}{115}{132}{Parametricity.PrivateRead.instance_source}
\par\penalty0

The instance target identity gives the entire ordered target list after applying the construction, with output order retained.
For the proof, constructor and boundary cases are separated explicitly; the defined maps and incidence equations are substituted.
\lean{Copying/ParametricityPrivateRead}{133}{133}{Parametricity.PrivateRead.instance_target}
\par\penalty0

The store source identity gives the entire ordered source list after applying the construction, rather than only its underlying set.
For the proof, constructor and boundary cases are separated explicitly; the defined maps and incidence equations are substituted.
\lean{Copying/ParametricityPrivateRead}{134}{134}{Parametricity.PrivateRead.store_source}
\par\penalty0

The store target identity gives the entire ordered target list after applying the construction, with output order retained.
For the proof, constructor and boundary cases are separated explicitly; the defined maps and incidence equations are substituted.
\lean{Copying/ParametricityPrivateRead}{135}{135}{Parametricity.PrivateRead.store_target}
\par\penalty0

Complete occurrence of the original dereference rule, with the real private store and the exact ordered three-contraction atom spine.
\lean{Copying/ParametricityPrivateRead}{139}{147}{Parametricity.PrivateRead.readOccurrence}
\par\penalty0

\subsection{Parametricity Read Prefix}
An actual dereference at one selected sharing input, with all other sharing inputs kept external. The prefix performs the original instance copy request and copy; it does not read the store early.

This defines the hole data for the construction just described. Boundary positions are kept distinct from newly inserted internal vertices.
\lean{Copying/ParametricityReadPrefix}{12}{12}{Parametricity.ReaderPrefix.hole}
\par\penalty0

This defines the label data for the construction just described. Boundary positions are kept distinct from newly inserted internal vertices.
\lean{Copying/ParametricityReadPrefix}{13}{15}{Parametricity.ReaderPrefix.label}
\par\penalty0

This defines the ordered source list for the construction just described. Boundary positions are kept distinct from newly inserted internal vertices.
\lean{Copying/ParametricityReadPrefix}{17}{20}{Parametricity.ReaderPrefix.source}
\par\penalty0

This defines the ordered target list for the construction just described. Boundary positions are kept distinct from newly inserted internal vertices.
\lean{Copying/ParametricityReadPrefix}{22}{25}{Parametricity.ReaderPrefix.target}
\par\penalty0

A real dereference edge precedes selected input zero. There are no identity edges at the remaining input positions.
\lean{Copying/ParametricityReadPrefix}{29}{100}{Parametricity.ReaderPrefix.frame}
\par\penalty0

The finite prefix at the original representation allowance.
\lean{Copying/ParametricityReadPrefix}{103}{104}{Parametricity.ReaderPrefix.pattern}
\par\penalty0

This defines the filling data for the construction just described. Boundary positions are kept distinct from newly inserted internal vertices.
\lean{Copying/ParametricityReadPrefix}{106}{110}{Parametricity.ReaderPrefix.filling}
\par\penalty0

This defines the finite typed net for the construction just described. Boundary positions are kept distinct from newly inserted internal vertices.
\lean{Copying/ParametricityReadPrefix}{112}{113}{Parametricity.ReaderPrefix.net}
\par\penalty0

This defines the active data for the construction just described. Boundary positions are kept distinct from newly inserted internal vertices.
\lean{Copying/ParametricityReadPrefix}{115}{116}{Parametricity.ReaderPrefix.active}
\par\penalty0

The inserted source's selected input is exactly the reader argument.
For the proof, constructor and boundary cases are separated explicitly; the defined maps and incidence equations are substituted.
\lean{Copying/ParametricityReadPrefix}{119}{121}{Parametricity.ReaderPrefix.selected_insert}
\par\penalty0

The search argument assertion is the compatibility law for the indicated component of this construction. It relates the actual finite coordinates or certificates, with the hypotheses stated in the linked declaration.
For the proof, constructor and boundary cases are separated explicitly; the defined maps and incidence equations are substituted.
\lean{Copying/ParametricityReadPrefix}{123}{129}{Parametricity.ReaderPrefix.search_argument}
\par\penalty0

The ready assertion is the compatibility law for the indicated component of this construction. It relates the actual finite coordinates or certificates, with the hypotheses stated in the linked declaration.
For the proof, constructor and boundary cases are separated explicitly; the defined maps and incidence equations are substituted.
\lean{Copying/ParametricityReadPrefix}{131}{136}{Parametricity.ReaderPrefix.ready}
\par\penalty0

Search in the unchanged selected component lifts through the literal reader frame, including the exact shared internal token coordinate.
For the proof, constructor and boundary cases are separated explicitly.
\lean{Copying/ParametricityReadPrefix}{140}{144}{Parametricity.ReaderPrefix.search_selected}
\par\penalty0

The rewrite selected assertion supplies the indicated actual machine transition or its transport through the common construction.
For the proof, constructor and boundary cases are separated explicitly; the defined maps and incidence equations are substituted.
\lean{Copying/ParametricityReadPrefix}{146}{152}{Parametricity.ReaderPrefix.rewrite_selected}
\par\penalty0

Five actual machine transitions expose the private read redex: enter its argument, request copying, copy the instance, return it, and refocus the reader. No store access or content copying is hidden in this prefix.
For the proof, the preceding component lemmas are composed or specialized at the stated parameters.
\lean{Copying/ParametricityReadPrefix}{157}{165}{Parametricity.ReaderPrefix.instance_prefix}
\par\penalty0

\subsection{Parametricity Read Tail}
The final three transitions of private readback are the literal read rewrite, the new distributor's copy request, and copying the stored numeral. The construction works for every atom tree and retains its other references.

This defines the numeral data for the construction just described. Boundary positions are kept distinct from newly inserted internal vertices.
\lean{Copying/ParametricityReadTail}{14}{14}{Parametricity.ReadTail.numeral}
\par\penalty0

This defines the before data for the construction just described. Boundary positions are kept distinct from newly inserted internal vertices.
\lean{Copying/ParametricityReadTail}{15}{15}{Parametricity.ReadTail.before}
\par\penalty0

This defines the after Read data for the construction just described. Boundary positions are kept distinct from newly inserted internal vertices.
\lean{Copying/ParametricityReadTail}{16}{16}{Parametricity.ReadTail.afterRead}
\par\penalty0

This defines the share data for the construction just described. Boundary positions are kept distinct from newly inserted internal vertices.
\lean{Copying/ParametricityReadTail}{18}{18}{Parametricity.ReadTail.share}
\par\penalty0

This defines the one data for the construction just described. Boundary positions are kept distinct from newly inserted internal vertices.
\lean{Copying/ParametricityReadTail}{19}{19}{Parametricity.ReadTail.one}
\par\penalty0

This defines the copy Edges data for the construction just described. Boundary positions are kept distinct from newly inserted internal vertices.
\lean{Copying/ParametricityReadTail}{21}{21}{Parametricity.ReadTail.copyEdges}
\par\penalty0

The copy Edges injective assertion states injectivity of the indicated coordinate map: equal image coordinates have equal original coordinates.
For the proof, the remaining natural-number index bounds follow by linear arithmetic; the defined maps and incidence equations are substituted.
\lean{Copying/ParametricityReadTail}{22}{32}{Parametricity.ReadTail.copyEdges_injective}
\par\penalty0

The copy label assertion verifies the indicated label compatibility or label exclusion at the original coordinate.
For the proof, constructor and boundary cases are separated explicitly; the remaining natural-number index bounds follow by linear arithmetic; the defined maps and incidence equations are substituted.
\lean{Copying/ParametricityReadTail}{34}{37}{Parametricity.ReadTail.copy_label}
\par\penalty0

This defines the copy Path data for the construction just described. Boundary positions are kept distinct from newly inserted internal vertices.
\lean{Copying/ParametricityReadTail}{39}{43}{Parametricity.ReadTail.copyPath}
\par\penalty0

This defines the copy Sides data for the construction just described. Boundary positions are kept distinct from newly inserted internal vertices.
\lean{Copying/ParametricityReadTail}{44}{45}{Parametricity.ReadTail.copySides}
\par\penalty0

The copy source identity gives the entire ordered source list after applying the construction, rather than only its underlying set.
For the proof, constructor and boundary cases are separated explicitly; the remaining natural-number index bounds follow by linear arithmetic; the defined maps and incidence equations are substituted.
\lean{Copying/ParametricityReadTail}{47}{51}{Parametricity.ReadTail.copy_source}
\par\penalty0

The copy target identity gives the entire ordered target list after applying the construction, with output order retained.
For the proof, constructor and boundary cases are separated explicitly; the remaining natural-number index bounds follow by linear arithmetic; the defined maps and incidence equations are substituted.
\lean{Copying/ParametricityReadTail}{53}{56}{Parametricity.ReadTail.copy_target}
\par\penalty0

The actual selected two-contraction spine, with unused store and weakening attachments exposed as boundary ports in their literal order.
\lean{Copying/ParametricityReadTail}{60}{64}{Parametricity.ReadTail.treeOccurrence}
\par\penalty0

The numeral Occurrence map is a complete occurrence, preserving original labels and every ordered source and target attachment.
\lean{Copying/ParametricityReadTail}{66}{67}{Parametricity.ReadTail.numeralOccurrence}
\par\penalty0

The copy boundary assertion is the compatibility law for the indicated component of this construction. It relates the actual finite coordinates or certificates, with the hypotheses stated in the linked declaration.
For the proof, constructor and boundary cases are separated explicitly; the remaining natural-number index bounds follow by linear arithmetic; the defined maps and incidence equations are substituted.
\lean{Copying/ParametricityReadTail}{69}{73}{Parametricity.ReadTail.copy_boundary}
\par\penalty0

The copy disjoint assertion excludes a shared edge or coordinate between the indicated independently inserted components.
For the proof, constructor and boundary cases are separated explicitly.
\lean{Copying/ParametricityReadTail}{75}{78}{Parametricity.ReadTail.copy_disjoint}
\par\penalty0

A complete match of the original copying rule, with literal stored 1.
\lean{Copying/ParametricityReadTail}{81}{84}{Parametricity.ReadTail.copyOccurrence}
\par\penalty0

The copy token assertion is the compatibility law for the indicated component of this construction. It relates the actual finite coordinates or certificates, with the hypotheses stated in the linked declaration.
For the proof, the defined maps and incidence equations are substituted.
\lean{Copying/ParametricityReadTail}{86}{91}{Parametricity.ReadTail.copy_token}
\par\penalty0

The actual finite cut replacement, retaining all unmatched atom-tree, store and weakening edges. This is the graph after the stored value is copied.
\lean{Copying/ParametricityReadTail}{95}{97}{Parametricity.ReadTail.afterCopy}
\par\penalty0

This defines the resume data for the construction just described. Boundary positions are kept distinct from newly inserted internal vertices.
\lean{Copying/ParametricityReadTail}{99}{101}{Parametricity.ReadTail.resume}
\par\penalty0

The copy rewrite assertion supplies the indicated actual machine transition or its transport through the common construction.
For the proof, the preceding component lemmas are composed or specialized at the stated parameters.
\lean{Copying/ParametricityReadTail}{103}{107}{Parametricity.ReadTail.copy_rewrite}
\par\penalty0

The read step assertion supplies the indicated actual machine transition or its transport through the common construction.
For the proof, the preceding component lemmas are composed or specialized at the stated parameters.
\lean{Copying/ParametricityReadTail}{109}{113}{Parametricity.ReadTail.read_step}
\par\penalty0

The request step assertion supplies the indicated actual machine transition or its transport through the common construction.
For the proof, the preceding component lemmas are composed or specialized at the stated parameters.
\lean{Copying/ParametricityReadTail}{115}{119}{Parametricity.ReadTail.request_step}
\par\penalty0

Three transitions, stopping with search at the freshly copied numeral. Its common passive return is deliberately not included in the nine-versus-one comparison of the paper.
For the proof, the preceding component lemmas are composed or specialized at the stated parameters.
\lean{Copying/ParametricityReadTail}{124}{128}{Parametricity.ReadTail.run}
\par\penalty0

\subsection{Parametricity Template}
Figure 20's complete literal pre-template, with independently arbitrary finite family sizes and arbitrary original term sharing tree. No robustness or observational equivalence is included as a premise or constructor.

The Literal data record the finite alternatives and their typed certificates used in this construction.
\lean{Copying/ParametricityTemplate}{12}{14}{Parametricity.Literal}
\par\penalty0

The Pair data record the finite alternatives and their typed certificates used in this construction.
\lean{Copying/ParametricityTemplate}{16}{18}{Parametricity.Pair}
\par\penalty0

The closed assertion excludes focus labels throughout the indicated graph, including recursively inside boxes.
For the proof, constructor and boundary cases are separated explicitly.
\lean{Copying/ParametricityTemplate}{20}{23}{Parametricity.Literal.closed}
\par\penalty0

The closed assertion excludes focus labels throughout the indicated graph, including recursively inside boxes.
For the proof, constructor and boundary cases are separated explicitly.
\lean{Copying/ParametricityTemplate}{25}{30}{Parametricity.Pair.closed}
\par\penalty0

No private term or atom wire escapes the complete family and its store.
For the proof, constructor and boundary cases are separated explicitly; the defined maps and incidence equations are substituted.
\lean{Copying/ParametricityTemplate}{33}{35}{Parametricity.Pair.noOutputs}
\par\penalty0

This defines the template data for the construction just described. Boundary positions are kept distinct from newly inserted internal vertices.
\lean{Copying/ParametricityTemplate}{37}{39}{Parametricity.template}
\par\penalty0

The literal assertion is the compatibility law for the indicated component of this construction. It relates the actual finite coordinates or certificates, with the hypotheses stated in the linked declaration.
For the proof, the preceding component lemmas are composed or specialized at the stated parameters.
\lean{Copying/ParametricityTemplate}{41}{43}{Parametricity.literal}
\par\penalty0

Output closure follows from the actual empty interface; it places no assumption on callback behavior or on the surrounding observing context.
For the proof, constructor and boundary cases are separated explicitly; the defined maps and incidence equations are substituted.
\lean{Copying/ParametricityTemplate}{47}{53}{Parametricity.template_outputClosed}
\par\penalty0

Every original entry returns in one search transition on both sides, including independently renamed graph carriers.
For the proof, constructor and boundary cases are separated explicitly.
\lean{Copying/ParametricityTemplate}{57}{65}{Parametricity.Pair.returns}
\par\penalty0

Input safety uses the same one-step return on both sides and relocates the shared token to backtracking. All original specimens and context classes are permitted, without any plugging-closure assumption.
For the proof, the defined maps and incidence equations are substituted.
\lean{Copying/ParametricityTemplate}{70}{95}{Parametricity.template_inputSafe}
\par\penalty0

